% Pertemuan 12 Praktikum Pemrograman (IF25-11002). Dihasilkan oleh tools/build.py dari tools/konten/p12.py.
% Kompilasi: xelatex P12.tex (dua kali). Catatan: xelatex -jobname=P12-catatan "\def\catatan{1}% Pertemuan 12 Praktikum Pemrograman (IF25-11002). Dihasilkan oleh tools/build.py dari tools/konten/p12.py.
% Kompilasi: xelatex P12.tex (dua kali). Catatan: xelatex -jobname=P12-catatan "\def\catatan{1}% Pertemuan 12 Praktikum Pemrograman (IF25-11002). Dihasilkan oleh tools/build.py dari tools/konten/p12.py.
% Kompilasi: xelatex P12.tex (dua kali). Catatan: xelatex -jobname=P12-catatan "\def\catatan{1}% Pertemuan 12 Praktikum Pemrograman (IF25-11002). Dihasilkan oleh tools/build.py dari tools/konten/p12.py.
% Kompilasi: xelatex P12.tex (dua kali). Catatan: xelatex -jobname=P12-catatan "\def\catatan{1}\input{P12.tex}"
\documentclass[t]{beamer}
\usetheme{ITERA}
\input{catatan-preamble.tex}
\renewcommand\ITERAmeeting{Pertemuan 12}
\begin{document}
\SlideJudul{IF25-11002 PRAKTIKUM PEMROGRAMAN}{Pertemuan 12\\Sorting}{Mengurutkan data dengan bubble sort dan selection sort, serta menelusuri setiap putarannya}{Tim Pengajar}{Tahun 2026. Sejajar dengan Pertemuan 12 Algoritma Pemrograman: Sorting.}
\note{Buka dengan menanyakan satu hal dari pertemuan lalu yang masih membingungkan.}

\begin{frame}{Dua mata kuliah, satu minggu}
\framesubtitle{Sebelum mulai}
Topik minggu ini sama di dua mata kuliah, dengan peran yang berbeda.
\vspace{10pt}

\begin{columns}[T]
\begin{column}{0.47\textwidth}
\begin{Blok}{Algoritma: Sorting}
Algoritma pengurutan gelembung dan seleksi, menelusuri isi larik per putaran, dan menghitung perbandingan serta pertukaran.
\end{Blok}
\end{column}
\begin{column}{0.47\textwidth}
\begin{BlokGelap}{Praktikum: Sorting}
Menulis pertukaran, bubble sort, dan selection sort di C++, termasuk mengurutkan dua array sejajar.
\end{BlokGelap}
\end{column}
\end{columns}
\note{Tanyakan apa yang sudah dibahas di Algoritma minggu ini. Gunakan jawabannya sebagai jembatan ke praktikum.}
\end{frame}

\begin{frame}{Saat pulang nanti, kalian bisa}
\framesubtitle{Target hari ini}
\begin{columns}[T]
\begin{column}{0.47\textwidth}
\begin{Langkah}{1}{Menerapkan}
Menulis program yang mengurutkan array naik maupun turun dengan bubble sort dan selection sort, termasuk mengurutkan dua array sejajar sekaligus

{\footnotesize\color{ITERAInkMuted}Sub-CPMK 12.1, CPMK0607}
\end{Langkah}
\end{column}
\begin{column}{0.47\textwidth}
\begin{Langkah}{2}{Menjelaskan}
Menelusuri isi array sesudah setiap putaran pengurutan dan menjelaskan banyak perbandingan serta pertukaran yang terjadi

{\footnotesize\color{ITERAInkMuted}Sub-CPMK 12.2, CPMK0608}
\end{Langkah}
\end{column}
\end{columns}
\vspace{8pt}
\begin{Catatan}
Dua kemampuan ini dinilai sama berat di Exit Ticket, tugas, UTS, dan UAS.
\end{Catatan}
\note{Bacakan target dengan pelan. Kata kuncinya menerapkan dan menjelaskan.}
\end{frame}

\begin{frame}{Ingat minggu lalu}
\framesubtitle{Review, 5 menit}
\begin{columns}[T]
\begin{column}{0.31\textwidth}
\begin{BlokGelap}{Pertanyaan 1}
Apa prasyarat binary search?
\end{BlokGelap}
\end{column}
\begin{column}{0.31\textwidth}
\begin{BlokGelap}{Pertanyaan 2}
Mengapa syarat binary search \texttt{low <= high}?
\end{BlokGelap}
\end{column}
\begin{column}{0.31\textwidth}
\begin{BlokGelap}{Pertanyaan 3}
Berapa langkah binary search untuk 1.000 data?
\end{BlokGelap}
\end{column}
\end{columns}
\vspace{8pt}
Jawab di kertas dulu, lalu cocokkan dengan teman sebelah.\note{Pilih dua mahasiswa untuk menjawab. Koreksi singkat, jangan mengajar ulang.}
\end{frame}

\Bagian{1}{Masalah pembuka}{Papan peringkat lomba coding}
\note{Masalah pembuka 10 menit.}

\begin{frame}[fragile]{Papan peringkat lomba coding}
\framesubtitle{Masalah minggu ini}
\begin{columns}[T]
\begin{column}{0.55\textwidth}
{\small Panitia lomba coding mencatat skor peserta sesuai urutan pendaftaran. Di akhir lomba, papan peringkat harus menampilkan peserta dari skor tertinggi.

\vspace{6pt}
Mencari skor tertinggi sudah bisa sejak Pertemuan 9. Tantangannya sekarang adalah menyusun semua peserta, dan nama harus tetap menempel pada skornya.\par}
\vspace{6pt}
\begin{Catatan}
\textbf{Diskusi 2 menit.} Bagaimana kalian mengurutkan lima kartu skor di atas meja dengan tangan? Langkah apa yang terus kalian ulang?
\end{Catatan}
\end{column}
\begin{column}{0.41\textwidth}
\begin{Keluaran}[]
1. Siti 95
2. Andi 88
3. Rani 82
4. Budi 75
5. Doni 70
\end{Keluaran}
\end{column}
\end{columns}
\note{Tampung beberapa jawaban tanpa menilai benar salah.}
\end{frame}

\begin{frame}{Dari langkah ke instruksi}
\framesubtitle{Sambungan dengan Algoritma}
\begin{columns}[T]
\begin{column}{0.31\textwidth}
\begin{Langkah}{1}{Pilih caranya}
Tukar pasangan bersebelahan yang terbalik, atau cari yang terbesar lalu taruh di depan.
\end{Langkah}
\end{column}
\begin{column}{0.31\textwidth}
\begin{Langkah}{2}{Ulangi per putaran}
Setiap putaran memastikan satu elemen sudah di tempat akhirnya.
\end{Langkah}
\end{column}
\begin{column}{0.31\textwidth}
\begin{Langkah}{3}{Terjemahkan}
Pertukaran tiga baris dan dua \texttt{for} bersarang.
\end{Langkah}
\end{column}
\end{columns}
\vspace{10pt}
Langkah 1 dan 2 dilatih di Algoritma. Langkah 3 kita kerjakan hari ini dengan C++.\note{Hubungkan eksplisit ke Algoritma minggu ini.}
\end{frame}

\Bagian{2}{Coba dulu, baru dijelaskan}{25 menit eksperimen di GDB Online}
\note{Struggle 25 menit. Berkeliling, jangan menjelaskan materi dulu.}

\begin{frame}{Misi kalian}
\framesubtitle{Struggle, 25 menit}
\begin{columns}[T]
\begin{column}{0.31\textwidth}
\begin{Langkah}{1}{Tukar}
Tukar isi \texttt{a} dan \texttt{b} sehingga \texttt{a = 7, b = 3} menjadi \texttt{a = 3, b = 7}.
\end{Langkah}
\end{column}
\begin{column}{0.31\textwidth}
\begin{Langkah}{2}{Satu putaran}
Pada \texttt{\{5, 1, 4, 2\}}, tukar setiap pasangan bersebelahan yang terbalik dari kiri ke kanan.
\end{Langkah}
\end{column}
\begin{column}{0.31\textwidth}
\begin{Langkah}{3}{Semua putaran}
Ulangi putaran itu sampai data terurut. Berapa putaran yang dibutuhkan?
\end{Langkah}
\end{column}
\end{columns}
\vspace{8pt}
\begin{columns}[T]
\begin{column}{0.47\textwidth}
\begin{Blok}{Boleh}
Bertanya ke teman, mengubah apa saja, sengaja membuat error.
\end{Blok}
\end{column}
\begin{column}{0.47\textwidth}
\begin{BlokAlert}{Belum dulu}
Mencari jawaban jadi di internet atau AI.
\end{BlokAlert}
\end{column}
\end{columns}
\note{Catat kesalahan yang sering muncul untuk dibahas di debrief.}
\end{frame}

\begin{frame}{Apa yang kalian temukan?}
\framesubtitle{Debrief struggle}
\begin{columns}[T]
\begin{column}{0.31\textwidth}
\begin{BlokGelap}{Pertanyaan 1}
Apa yang terjadi bila menukar dengan \texttt{a = b; b = a;}?
\end{BlokGelap}
\end{column}
\begin{column}{0.31\textwidth}
\begin{BlokGelap}{Pertanyaan 2}
Sesudah satu putaran, elemen mana yang pasti sudah di tempatnya?
\end{BlokGelap}
\end{column}
\begin{column}{0.31\textwidth}
\begin{BlokGelap}{Pertanyaan 3}
Apakah selalu butuh n − 1 putaran?
\end{BlokGelap}
\end{column}
\end{columns}
\vspace{10pt}
Jawaban kalian akan kita uji di bagian berikutnya.\note{Tulis jawaban di papan; buktikan atau patahkan di bagian materi.}
\end{frame}

\Bagian{3}{Mengurutkan langkah demi langkah}{Pertukaran, bubble sort, selection sort, dan array sejajar}
\note{Materi 40 menit. Tunjukkan setiap contoh langsung di GDB Online.}

\begin{frame}[fragile]{Menukar dua nilai}
\framesubtitle{Variabel sementara}
\begin{columns}[T]
\begin{column}{0.55\textwidth}
\begin{Kode}[]{tukar.cpp}
int a = 7, b = 3;
int t = a;
a = b;
b = t;
cout << a << " " << b << endl;
\end{Kode}
\end{column}
\begin{column}{0.41\textwidth}
\only<1>{\begin{TebakDulu}
{\ITERAKickerFont\fontsize{10pt}{12pt}\selectfont\color{ITERAAlert}TEBAK DULU}\par\vspace{4pt}
Sebelum dijalankan, tulis keluarannya di kertas.
\end{TebakDulu}
}\only<2>{\KeluaranFile[]{keluaran/P12/k001.txt}
\begin{Catatan}
Tanpa \texttt{t}, nilai lama \texttt{a} hilang begitu ditimpa.
\end{Catatan}
}
\end{column}
\end{columns}
\note<1>{Minta mahasiswa menebak keluaran sebelum membuka slide berikutnya. Tanyakan satu atau dua tebakan yang berbeda, lalu buktikan.}
\end{frame}

\begin{frame}{Bubble sort: satu putaran}
\framesubtitle{Pasangan bersebelahan}
Putaran pertama pada \{5, 1, 4, 2, 8\}: bandingkan setiap pasangan, tukar bila terbalik.
\vspace{8pt}

\small
\begin{tabular}{@{}>{\raggedright\arraybackslash}p{172pt}>{\raggedright\arraybackslash}p{326pt}>{\raggedright\arraybackslash}p{364pt}@{}}
\kepala{Pasangan} & \kepala{Isi array sesudahnya} & \kepala{Keterangan}\\
a[0], a[1] & 1 5 4 2 8 & 5 > 1, tukar\\
\barisgenap a[1], a[2] & 1 4 5 2 8 & 5 > 4, tukar\\
a[2], a[3] & 1 4 2 5 8 & 5 > 2, tukar\\
\barisgenap a[3], a[4] & 1 4 2 5 8 & 5 < 8, tetap\\
\garisdasar
\end{tabular}
\vspace{10pt}

Sesudah putaran pertama, elemen terbesar pasti sudah di ujung kanan, seperti gelembung yang naik ke permukaan.
\end{frame}

\begin{frame}[fragile]{Bubble sort di C++}
\framesubtitle{Dua perulangan bersarang}
\begin{columns}[T]
\begin{column}{0.60\textwidth}
\begin{Kode}[listing options={style=ITERACodeKecil}]{bubble.cpp}
int a[5] = {5, 1, 4, 2, 8};
int n = 5;
for (int i = 0; i < n - 1; i++) {
    for (int j = 0; j < n - 1 - i; j++) {
        if (a[j] > a[j + 1]) {
            int t = a[j]; a[j] = a[j + 1]; a[j + 1] = t;
        }
    }
}
for (int k = 0; k < n; k++) cout << a[k] << " ";
cout << endl;
\end{Kode}
\end{column}
\begin{column}{0.36\textwidth}
\only<1>{\begin{TebakDulu}
{\ITERAKickerFont\fontsize{10pt}{12pt}\selectfont\color{ITERAAlert}TEBAK DULU}\par\vspace{4pt}
Sebelum dijalankan, tulis keluarannya di kertas.
\end{TebakDulu}
}\only<2>{\KeluaranFile[listing options={style=ITERAPlainKecil}]{keluaran/P12/k002.txt}
\begin{Catatan}
Batas dalam \texttt{n - 1 - i}: setiap putaran, satu elemen di ujung sudah final. \texttt{j + 1} tidak boleh melewati n − 1.
\end{Catatan}
}
\end{column}
\end{columns}
\note<1>{Minta mahasiswa menebak keluaran sebelum membuka slide berikutnya. Tanyakan satu atau dua tebakan yang berbeda, lalu buktikan.}
\end{frame}

\begin{frame}[fragile]{Selection sort}
\framesubtitle{Pilih yang terkecil}
\begin{columns}[T]
\begin{column}{0.60\textwidth}
\begin{Kode}[listing options={style=ITERACodeKecil}]{selection.cpp}
int a[5] = {29, 10, 14, 37, 13};
for (int i = 0; i < 4; i++) {
    int iMin = i;
    for (int j = i + 1; j < 5; j++) {
        if (a[j] < a[iMin]) iMin = j;
    }
    int t = a[i]; a[i] = a[iMin]; a[iMin] = t;
}
for (int k = 0; k < 5; k++) cout << a[k] << " ";
cout << endl;
\end{Kode}
\end{column}
\begin{column}{0.36\textwidth}
\only<1>{\begin{TebakDulu}
{\ITERAKickerFont\fontsize{10pt}{12pt}\selectfont\color{ITERAAlert}TEBAK DULU}\par\vspace{4pt}
Sebelum dijalankan, tulis keluarannya di kertas.
\end{TebakDulu}
}\only<2>{\KeluaranFile[listing options={style=ITERAPlainKecil}]{keluaran/P12/k003.txt}
\begin{Catatan}
Paling banyak satu pertukaran per putaran. Untuk urutan turun, cari yang terbesar.
\end{Catatan}
}
\end{column}
\end{columns}
\note<1>{Minta mahasiswa menebak keluaran sebelum membuka slide berikutnya. Tanyakan satu atau dua tebakan yang berbeda, lalu buktikan.}
\end{frame}

\begin{frame}[fragile]{Mengurutkan dua array sejajar}
\framesubtitle{Nama ikut skornya}
\begin{columns}[T]
\begin{column}{0.60\textwidth}
\begin{Kode}[listing options={style=ITERACodeKecil}]{sejajar.cpp}
string nama[4] = {"Rani", "Budi", "Siti", "Andi"};
int skor[4] = {82, 75, 95, 88};
for (int i = 0; i < 3; i++) {
    int iMaks = i;
    for (int j = i + 1; j < 4; j++)
        if (skor[j] > skor[iMaks]) iMaks = j;
    swap(skor[i], skor[iMaks]);
    swap(nama[i], nama[iMaks]);
}
for (int i = 0; i < 4; i++)
    cout << nama[i] << " " << skor[i] << endl;
\end{Kode}
\end{column}
\begin{column}{0.36\textwidth}
\only<1>{\begin{TebakDulu}
{\ITERAKickerFont\fontsize{10pt}{12pt}\selectfont\color{ITERAAlert}TEBAK DULU}\par\vspace{4pt}
Sebelum dijalankan, tulis keluarannya di kertas.
\end{TebakDulu}
}\only<2>{\KeluaranFile[listing options={style=ITERAPlainKecil}]{keluaran/P12/k004.txt}
\begin{Catatan}
Setiap pertukaran pada array kunci harus diikuti pertukaran yang sama pada array pasangannya. \texttt{swap} dari pustaka standar menukar dua nilai.
\end{Catatan}
}
\end{column}
\end{columns}
\note<1>{Minta mahasiswa menebak keluaran sebelum membuka slide berikutnya. Tanyakan satu atau dua tebakan yang berbeda, lalu buktikan.}
\end{frame}

\begin{frame}[fragile]{Tebak keluarannya}
\framesubtitle{Latihan menjelaskan, 3 menit}
\begin{columns}[T]
\begin{column}{0.56\textwidth}
\begin{Kode}[listing options={style=ITERACodeKecil}]{tebak.cpp}
int a[4] = {4, 3, 2, 1};
int tukar = 0;
for (int i = 0; i < 1; i++) {
    for (int j = 0; j < 3 - i; j++) {
        if (a[j] > a[j + 1]) {
            swap(a[j], a[j + 1]);
            tukar++;
        }
    }
}
for (int k = 0; k < 4; k++) cout << a[k];
cout << " " << tukar << endl;
\end{Kode}
\end{column}
\begin{column}{0.40\textwidth}
\begin{Catatan}
Tulis keluarannya di kertas, karakter demi karakter. Jangan dijalankan dulu.
\end{Catatan}
\end{column}
\end{columns}
\note{Contoh soal Bagian B. Beri waktu 3 menit sebelum slide jawaban.}
\end{frame}

\begin{frame}[fragile]{Tebak keluarannya: jawaban}
\framesubtitle{Latihan menjelaskan}
\begin{columns}[T]
\begin{column}{0.56\textwidth}
\begin{Kode}[listing options={style=ITERACodeKecil}]{tebak.cpp}
int a[4] = {4, 3, 2, 1};
int tukar = 0;
for (int i = 0; i < 1; i++) {
    for (int j = 0; j < 3 - i; j++) {
        if (a[j] > a[j + 1]) {
            swap(a[j], a[j + 1]);
            tukar++;
        }
    }
}
for (int k = 0; k < 4; k++) cout << a[k];
cout << " " << tukar << endl;
\end{Kode}
\end{column}
\begin{column}{0.40\textwidth}
\begin{Keluaran}[]
3214 3
\end{Keluaran}
\begin{Catatan}
Hanya satu putaran (i = 0). Angka 4 terdorong ke ujung lewat tiga pertukaran, tetapi sisanya belum terurut: 3 2 1 4.
\end{Catatan}
\end{column}
\end{columns}
\note{Tanyakan jawaban yang paling sering salah, lalu telusuri bersama.}
\end{frame}

\begin{frame}{Error yang sering muncul minggu ini}
\framesubtitle{Pesan diambil langsung dari g++}
\footnotesize
\begin{tabular}{@{}>{\raggedright\arraybackslash}p{208pt}>{\raggedright\arraybackslash}p{256pt}>{\raggedright\arraybackslash}p{398pt}@{}}
\kepala{Kesalahan} & \kepala{Contoh} & \kepala{Pesan pertama dari g++}\\
swap tanpa nilai yang bisa diubah & \texttt{swap(a[0], 5)} & \texttt{error: no matching function for call to 'swap(int\&, int)'}\\
\barisgenap Tukar string dengan int & \texttt{swap(nama[i], skor[i])} & \texttt{error: no matching function for call to 'swap(std::string\&, int\&)'}\\
Variabel sementara tak dipakai & \texttt{int t; tidak dipakai} & \texttt{warning: unused variable 't'}\\
\barisgenap Indeks bukan bilangan bulat & \texttt{a[j + 0.5]} & \texttt{error: invalid types 'int [3][double]' for array subscript}\\
Ukuran array tidak cocok & \texttt{int a[3] = \{1, 2, 3, 4\}} & \texttt{error: too many initializers for 'int [3]'}\\
\garisdasar
\end{tabular}
\vspace{10pt}

{\small Pesan di tabel ini dihasilkan dengan g++ dan opsi -Wall, sama seperti di cek.sh.}
\note{Minta mahasiswa memotret slide ini.}
\end{frame}

\Bagian{4}{Hands-on bertingkat}{45 menit, dari dasar sampai tantangan}
\note{Mahasiswa membuka folder 02 Hands-on/P12 - Sorting - Hands-on.}

\begin{frame}{Peta hands-on}
\framesubtitle{Folder 02 Hands-on/P12 - Sorting - Hands-on}
\small
\begin{tabular}{@{}>{\raggedright\arraybackslash}p{36pt}>{\raggedright\arraybackslash}p{267pt}>{\raggedright\arraybackslash}p{110pt}>{\raggedright\arraybackslash}p{83pt}>{\raggedright\arraybackslash}p{341pt}@{}}
\kepala{No} & \kepala{File} & \kepala{Tingkat} & \kepala{Waktu} & \kepala{Yang dilatih}\\
1 & \texttt{handson1-bubble.cpp} & Dasar & 10' & bubble sort naik dengan jejak per putaran\\
\barisgenap 2 & \texttt{handson2-peringkat.cpp} & Menengah & 15' & selection sort turun, dua array sejajar\\
3 & \texttt{handson3-tukar.cpp} & Tantangan & 20' & telusuri pertukaran, perbaiki tiga kesalahan\\
\barisgenap + & \texttt{bonus-hitung.cpp} & Pengayaan & sisa & bubble sort berhenti dini, hitung kerja\\
\garisdasar
\end{tabular}
\vspace{10pt}

Kerjakan berurutan. Cocokkan keluaran dengan folder \texttt{expected/}; latihan yang membaca masukan memakai data di folder \texttt{input/}.\note{Saat berkeliling, minta mahasiswa yang mengerjakan latihan tantangan menjelaskan blok PENJELASAN mereka.}
\end{frame}

\begin{frame}{Cara mengecek dan aturannya}
\framesubtitle{Hands-on}
\begin{columns}[T]
\begin{column}{0.47\textwidth}
\begin{Blok}{Di GDB Online}
Salin isi file latihan, lengkapi TODO, klik Run. Bila program membaca masukan, ketik isi file di \texttt{input/}. Bandingkan dengan \texttt{expected/}.
\end{Blok}
\end{column}
\begin{column}{0.47\textwidth}
\begin{BlokGelap}{Di komputer sendiri}
Jalankan \texttt{bash cek.sh}. Skrip mengompilasi dengan \texttt{-Wall -Werror}, memberi masukan dari \texttt{input/}, lalu mencocokkan keluaran.
\end{BlokGelap}
\end{column}
\end{columns}
\vspace{6pt}
\begin{Catatan}
AI boleh untuk memahami pesan error, tidak untuk meminta solusi utuh. Exit Ticket dikerjakan tanpa bantuan apa pun.
\end{Catatan}
\note{Hands-on tidak dinilai langsung; yang dinilai Exit Ticket dan tugas.}
\end{frame}

\Bagian{5}{Exit Ticket dan tugas}{25 menit terakhir, lalu pekerjaan rumah}
\note{Mulai Exit Ticket tepat waktu.}

\begin{frame}{Exit Ticket 12}
\framesubtitle{Individual, 25 menit, tanpa AI dan internet}
\small
\begin{tabular}{@{}>{\raggedright\arraybackslash}p{60pt}>{\raggedright\arraybackslash}p{560pt}>{\raggedright\arraybackslash}p{80pt}>{\raggedright\arraybackslash}p{150pt}@{}}
\kepala{Soal} & \kepala{Isi} & \kepala{Poin} & \kepala{CPMK}\\
A1 & Tulis bubble sort yang mengurutkan array turun & 25 & CPMK0607\\
\barisgenap A2 & Urutkan dua array sejajar (nama dan nilai) berdasarkan nilai & 25 & CPMK0607\\
B1 & Tuliskan isi array sesudah setiap putaran selection sort & 20 & CPMK0608\\
\barisgenap B2 & Jelaskan kesalahan pada potongan kode pertukaran & 20 & CPMK0608\\
B3 & Bandingkan banyak pertukaran bubble dan selection sort untuk data yang diberikan & 10 & CPMK0608\\
\garisdasar
\end{tabular}
\vspace{10pt}

{\small Soal A dinilai dari keluaran yang tepat sama. Soal B dinilai dari ketepatan dan alasan. Pelanggaran aturan membuat nilai Exit Ticket ini 0.}
\note{Soal lengkap dibagikan terpisah dan tidak disimpan di repo publik.}
\end{frame}

\begin{frame}{Tugas 12: Sistem Leaderboard}
\framesubtitle{Dikumpulkan sebelum Pertemuan 13}
\begin{columns}[T]
\begin{column}{0.47\textwidth}
\begin{Blok}{Bagian A, program, 50 poin}
Buat program leaderboard dengan menu berikut untuk maksimal 10 data (nama satu kata dan nilai). Ketentuannya di slide berikut.
\end{Blok}
\end{column}
\begin{column}{0.47\textwidth}
\begin{BlokGelap}{Bagian B, penjelasan, 50 poin}
Komentar maksimal 150 kata dan tabel isi array sesudah setiap putaran selection sort untuk lima data contoh. Jelaskan mengapa nama harus ikut ditukar, dan mengapa binary search hanya dipakai sesudah pengurutan abjad.
\end{BlokGelap}
\end{column}
\end{columns}
\vspace{8pt}
{\small Data di tugas adalah data kalian sendiri. Rubrik lengkap ada di Modul Belajar 12.}
\note{Ingatkan bahwa penjelasan disalin dari AI mudah dikenali karena tidak menyebut pengalaman mereka sendiri.}
\end{frame}

\begin{frame}{Ketentuan Tugas 12}
\framesubtitle{Sistem Leaderboard}
\footnotesize
\begin{tabular}{@{}>{\raggedright\arraybackslash}p{87pt}>{\raggedright\arraybackslash}p{788pt}@{}}
\kepala{No} & \kepala{Ketentuan Bagian A}\\
1 & Tambah data peserta dengan validasi nilai 0 sampai 100\\
\barisgenap 2 & Urutkan berdasarkan nilai turun dengan selection sort, lalu tampilkan leaderboard\\
3 & Urutkan berdasarkan nama (abjad) dengan bubble sort\\
\barisgenap 4 & Cari nama dengan linear search; tampilkan peringkat dan nilainya\\
5 & Cari nama dengan binary search pada data yang sudah urut abjad\\
\barisgenap 6 & Tampilkan banyak perbandingan dan pertukaran setiap algoritma yang dipakai, lalu keluar\\
Bonus & Bandingkan banyak langkah linear dan binary search untuk data yang sama, dan tampilkan dalam satu tabel.\\
\garisdasar
\end{tabular}
\note{Bacakan ketentuan satu per satu dan tanyakan apakah ada yang belum jelas.}
\end{frame}

\begin{frame}{Rubrik Tugas 12}
\framesubtitle{Empat level, sama di semua instrumen}
\footnotesize
\begin{tabular}{@{}>{\raggedright\arraybackslash}p{166pt}>{\raggedright\arraybackslash}p{44pt}>{\raggedright\arraybackslash}p{166pt}>{\raggedright\arraybackslash}p{156pt}>{\raggedright\arraybackslash}p{146pt}>{\raggedright\arraybackslash}p{146pt}@{}}
\kepala{Kriteria} & \kepala{Poin} & \kepala{Sangat baik 85--100} & \kepala{Baik 70--84} & \kepala{Cukup 55--69} & \kepala{Kurang <55}\\
A1 Program berjalan & 20 & tanpa error dan peringatan & ada peringatan & perlu satu perbaikan & tidak terkompilasi\\
\barisgenap A2 Kelengkapan menu & 15 & dua sort, dua search, penghitung & satu fitur kurang & dua kurang & tiga atau lebih kurang\\
A3 Ketepatan urutan dan pencarian & 15 & semua tepat termasuk nilai sama & satu kasus salah & dua salah & sebagian besar salah\\
\barisgenap B1 Jejak putaran dan alasan & 30 & tabel tepat dan alasan jelas & satu putaran keliru & tanpa alasan & tidak ada atau disalin\\
B2 Cerita error dan pengujian & 20 & pesan, sebab, perbaikan, uji & pesan tidak dikutip & hanya perbaikan & tidak ada\\
\garisdasar
\end{tabular}
\vspace{8pt}

{\small Nilai kriteria = poin × skor level. Bagian A total 50, Bagian B total 50.}
\note{Rubrik ini sama strukturnya setiap minggu; yang berubah hanya isi kriteria A2, A3, dan B1.}
\end{frame}

\begin{frame}{Sebelum Pertemuan 13}
\framesubtitle{Persiapan}
\begin{columns}[T]
\begin{column}{0.31\textwidth}
\begin{Langkah}{1}{Selesaikan}
Hands-on yang belum lulus \texttt{cek.sh}, terutama latihan tantangan.
\end{Langkah}
\end{column}
\begin{column}{0.31\textwidth}
\begin{Langkah}{2}{Kumpulkan}
Tugas 12 sebelum Pertemuan 13 dimulai.
\end{Langkah}
\end{column}
\begin{column}{0.31\textwidth}
\begin{Langkah}{3}{Baca}
Modul Belajar 12 untuk mengulang, lalu bagian awal modul berikutnya.
\end{Langkah}
\end{column}
\end{columns}
\vspace{10pt}
Pertemuan 13 membahas fungsi: memecah program menjadi bagian bernama yang bisa dipanggil berulang, termasuk fungsi pengurutan dan pencarian.\note{Tanyakan satu hal yang paling membingungkan hari ini untuk dibuka di review minggu depan.}
\end{frame}

\SlidePenutup{Terima kasih}{Sampai jumpa di Pertemuan 13: fungsi dan modularitas}
\note{Slide penutup.}
\end{document}
"
\documentclass[t]{beamer}
\usetheme{ITERA}
% Mode catatan pembicara: aktif bila \catatan didefinisikan sebelum \input.
\ifdefined\catatan
  \setbeameroption{show only notes}
  \setbeamertemplate{note page}{\vspace*{20pt}\hspace*{4pt}\begin{minipage}{900pt}
    {\ITERAKickerFont\fontsize{11pt}{14pt}\selectfont\color{ITERAGoldText}CATATAN PEMBICARA \ \ |\ \ SLIDE \insertframenumber}\par\vspace{4pt}
    {\ITERADisplay\bfseries\fontsize{22pt}{28pt}\selectfont\color{ITERAStructure}\insertframetitle}\par\vspace{12pt}
    \fontsize{15pt}{23pt}\selectfont\insertnote\end{minipage}}
\fi
\tikzset{fase/.style={draw=none,fill=#1,rounded corners=3pt,minimum width=158pt,minimum height=118pt,text width=146pt,align=center}}

\renewcommand\ITERAmeeting{Pertemuan 12}
\begin{document}
\SlideJudul{IF25-11002 PRAKTIKUM PEMROGRAMAN}{Pertemuan 12\\Sorting}{Mengurutkan data dengan bubble sort dan selection sort, serta menelusuri setiap putarannya}{Tim Pengajar}{Tahun 2026. Sejajar dengan Pertemuan 12 Algoritma Pemrograman: Sorting.}
\note{Buka dengan menanyakan satu hal dari pertemuan lalu yang masih membingungkan.}

\begin{frame}{Dua mata kuliah, satu minggu}
\framesubtitle{Sebelum mulai}
Topik minggu ini sama di dua mata kuliah, dengan peran yang berbeda.
\vspace{10pt}

\begin{columns}[T]
\begin{column}{0.47\textwidth}
\begin{Blok}{Algoritma: Sorting}
Algoritma pengurutan gelembung dan seleksi, menelusuri isi larik per putaran, dan menghitung perbandingan serta pertukaran.
\end{Blok}
\end{column}
\begin{column}{0.47\textwidth}
\begin{BlokGelap}{Praktikum: Sorting}
Menulis pertukaran, bubble sort, dan selection sort di C++, termasuk mengurutkan dua array sejajar.
\end{BlokGelap}
\end{column}
\end{columns}
\note{Tanyakan apa yang sudah dibahas di Algoritma minggu ini. Gunakan jawabannya sebagai jembatan ke praktikum.}
\end{frame}

\begin{frame}{Saat pulang nanti, kalian bisa}
\framesubtitle{Target hari ini}
\begin{columns}[T]
\begin{column}{0.47\textwidth}
\begin{Langkah}{1}{Menerapkan}
Menulis program yang mengurutkan array naik maupun turun dengan bubble sort dan selection sort, termasuk mengurutkan dua array sejajar sekaligus

{\footnotesize\color{ITERAInkMuted}Sub-CPMK 12.1, CPMK0607}
\end{Langkah}
\end{column}
\begin{column}{0.47\textwidth}
\begin{Langkah}{2}{Menjelaskan}
Menelusuri isi array sesudah setiap putaran pengurutan dan menjelaskan banyak perbandingan serta pertukaran yang terjadi

{\footnotesize\color{ITERAInkMuted}Sub-CPMK 12.2, CPMK0608}
\end{Langkah}
\end{column}
\end{columns}
\vspace{8pt}
\begin{Catatan}
Dua kemampuan ini dinilai sama berat di Exit Ticket, tugas, UTS, dan UAS.
\end{Catatan}
\note{Bacakan target dengan pelan. Kata kuncinya menerapkan dan menjelaskan.}
\end{frame}

\begin{frame}{Ingat minggu lalu}
\framesubtitle{Review, 5 menit}
\begin{columns}[T]
\begin{column}{0.31\textwidth}
\begin{BlokGelap}{Pertanyaan 1}
Apa prasyarat binary search?
\end{BlokGelap}
\end{column}
\begin{column}{0.31\textwidth}
\begin{BlokGelap}{Pertanyaan 2}
Mengapa syarat binary search \texttt{low <= high}?
\end{BlokGelap}
\end{column}
\begin{column}{0.31\textwidth}
\begin{BlokGelap}{Pertanyaan 3}
Berapa langkah binary search untuk 1.000 data?
\end{BlokGelap}
\end{column}
\end{columns}
\vspace{8pt}
Jawab di kertas dulu, lalu cocokkan dengan teman sebelah.\note{Pilih dua mahasiswa untuk menjawab. Koreksi singkat, jangan mengajar ulang.}
\end{frame}

\Bagian{1}{Masalah pembuka}{Papan peringkat lomba coding}
\note{Masalah pembuka 10 menit.}

\begin{frame}[fragile]{Papan peringkat lomba coding}
\framesubtitle{Masalah minggu ini}
\begin{columns}[T]
\begin{column}{0.55\textwidth}
{\small Panitia lomba coding mencatat skor peserta sesuai urutan pendaftaran. Di akhir lomba, papan peringkat harus menampilkan peserta dari skor tertinggi.

\vspace{6pt}
Mencari skor tertinggi sudah bisa sejak Pertemuan 9. Tantangannya sekarang adalah menyusun semua peserta, dan nama harus tetap menempel pada skornya.\par}
\vspace{6pt}
\begin{Catatan}
\textbf{Diskusi 2 menit.} Bagaimana kalian mengurutkan lima kartu skor di atas meja dengan tangan? Langkah apa yang terus kalian ulang?
\end{Catatan}
\end{column}
\begin{column}{0.41\textwidth}
\begin{Keluaran}[]
1. Siti 95
2. Andi 88
3. Rani 82
4. Budi 75
5. Doni 70
\end{Keluaran}
\end{column}
\end{columns}
\note{Tampung beberapa jawaban tanpa menilai benar salah.}
\end{frame}

\begin{frame}{Dari langkah ke instruksi}
\framesubtitle{Sambungan dengan Algoritma}
\begin{columns}[T]
\begin{column}{0.31\textwidth}
\begin{Langkah}{1}{Pilih caranya}
Tukar pasangan bersebelahan yang terbalik, atau cari yang terbesar lalu taruh di depan.
\end{Langkah}
\end{column}
\begin{column}{0.31\textwidth}
\begin{Langkah}{2}{Ulangi per putaran}
Setiap putaran memastikan satu elemen sudah di tempat akhirnya.
\end{Langkah}
\end{column}
\begin{column}{0.31\textwidth}
\begin{Langkah}{3}{Terjemahkan}
Pertukaran tiga baris dan dua \texttt{for} bersarang.
\end{Langkah}
\end{column}
\end{columns}
\vspace{10pt}
Langkah 1 dan 2 dilatih di Algoritma. Langkah 3 kita kerjakan hari ini dengan C++.\note{Hubungkan eksplisit ke Algoritma minggu ini.}
\end{frame}

\Bagian{2}{Coba dulu, baru dijelaskan}{25 menit eksperimen di GDB Online}
\note{Struggle 25 menit. Berkeliling, jangan menjelaskan materi dulu.}

\begin{frame}{Misi kalian}
\framesubtitle{Struggle, 25 menit}
\begin{columns}[T]
\begin{column}{0.31\textwidth}
\begin{Langkah}{1}{Tukar}
Tukar isi \texttt{a} dan \texttt{b} sehingga \texttt{a = 7, b = 3} menjadi \texttt{a = 3, b = 7}.
\end{Langkah}
\end{column}
\begin{column}{0.31\textwidth}
\begin{Langkah}{2}{Satu putaran}
Pada \texttt{\{5, 1, 4, 2\}}, tukar setiap pasangan bersebelahan yang terbalik dari kiri ke kanan.
\end{Langkah}
\end{column}
\begin{column}{0.31\textwidth}
\begin{Langkah}{3}{Semua putaran}
Ulangi putaran itu sampai data terurut. Berapa putaran yang dibutuhkan?
\end{Langkah}
\end{column}
\end{columns}
\vspace{8pt}
\begin{columns}[T]
\begin{column}{0.47\textwidth}
\begin{Blok}{Boleh}
Bertanya ke teman, mengubah apa saja, sengaja membuat error.
\end{Blok}
\end{column}
\begin{column}{0.47\textwidth}
\begin{BlokAlert}{Belum dulu}
Mencari jawaban jadi di internet atau AI.
\end{BlokAlert}
\end{column}
\end{columns}
\note{Catat kesalahan yang sering muncul untuk dibahas di debrief.}
\end{frame}

\begin{frame}{Apa yang kalian temukan?}
\framesubtitle{Debrief struggle}
\begin{columns}[T]
\begin{column}{0.31\textwidth}
\begin{BlokGelap}{Pertanyaan 1}
Apa yang terjadi bila menukar dengan \texttt{a = b; b = a;}?
\end{BlokGelap}
\end{column}
\begin{column}{0.31\textwidth}
\begin{BlokGelap}{Pertanyaan 2}
Sesudah satu putaran, elemen mana yang pasti sudah di tempatnya?
\end{BlokGelap}
\end{column}
\begin{column}{0.31\textwidth}
\begin{BlokGelap}{Pertanyaan 3}
Apakah selalu butuh n − 1 putaran?
\end{BlokGelap}
\end{column}
\end{columns}
\vspace{10pt}
Jawaban kalian akan kita uji di bagian berikutnya.\note{Tulis jawaban di papan; buktikan atau patahkan di bagian materi.}
\end{frame}

\Bagian{3}{Mengurutkan langkah demi langkah}{Pertukaran, bubble sort, selection sort, dan array sejajar}
\note{Materi 40 menit. Tunjukkan setiap contoh langsung di GDB Online.}

\begin{frame}[fragile]{Menukar dua nilai}
\framesubtitle{Variabel sementara}
\begin{columns}[T]
\begin{column}{0.55\textwidth}
\begin{Kode}[]{tukar.cpp}
int a = 7, b = 3;
int t = a;
a = b;
b = t;
cout << a << " " << b << endl;
\end{Kode}
\end{column}
\begin{column}{0.41\textwidth}
\only<1>{\begin{TebakDulu}
{\ITERAKickerFont\fontsize{10pt}{12pt}\selectfont\color{ITERAAlert}TEBAK DULU}\par\vspace{4pt}
Sebelum dijalankan, tulis keluarannya di kertas.
\end{TebakDulu}
}\only<2>{\KeluaranFile[]{keluaran/P12/k001.txt}
\begin{Catatan}
Tanpa \texttt{t}, nilai lama \texttt{a} hilang begitu ditimpa.
\end{Catatan}
}
\end{column}
\end{columns}
\note<1>{Minta mahasiswa menebak keluaran sebelum membuka slide berikutnya. Tanyakan satu atau dua tebakan yang berbeda, lalu buktikan.}
\end{frame}

\begin{frame}{Bubble sort: satu putaran}
\framesubtitle{Pasangan bersebelahan}
Putaran pertama pada \{5, 1, 4, 2, 8\}: bandingkan setiap pasangan, tukar bila terbalik.
\vspace{8pt}

\small
\begin{tabular}{@{}>{\raggedright\arraybackslash}p{172pt}>{\raggedright\arraybackslash}p{326pt}>{\raggedright\arraybackslash}p{364pt}@{}}
\kepala{Pasangan} & \kepala{Isi array sesudahnya} & \kepala{Keterangan}\\
a[0], a[1] & 1 5 4 2 8 & 5 > 1, tukar\\
\barisgenap a[1], a[2] & 1 4 5 2 8 & 5 > 4, tukar\\
a[2], a[3] & 1 4 2 5 8 & 5 > 2, tukar\\
\barisgenap a[3], a[4] & 1 4 2 5 8 & 5 < 8, tetap\\
\garisdasar
\end{tabular}
\vspace{10pt}

Sesudah putaran pertama, elemen terbesar pasti sudah di ujung kanan, seperti gelembung yang naik ke permukaan.
\end{frame}

\begin{frame}[fragile]{Bubble sort di C++}
\framesubtitle{Dua perulangan bersarang}
\begin{columns}[T]
\begin{column}{0.60\textwidth}
\begin{Kode}[listing options={style=ITERACodeKecil}]{bubble.cpp}
int a[5] = {5, 1, 4, 2, 8};
int n = 5;
for (int i = 0; i < n - 1; i++) {
    for (int j = 0; j < n - 1 - i; j++) {
        if (a[j] > a[j + 1]) {
            int t = a[j]; a[j] = a[j + 1]; a[j + 1] = t;
        }
    }
}
for (int k = 0; k < n; k++) cout << a[k] << " ";
cout << endl;
\end{Kode}
\end{column}
\begin{column}{0.36\textwidth}
\only<1>{\begin{TebakDulu}
{\ITERAKickerFont\fontsize{10pt}{12pt}\selectfont\color{ITERAAlert}TEBAK DULU}\par\vspace{4pt}
Sebelum dijalankan, tulis keluarannya di kertas.
\end{TebakDulu}
}\only<2>{\KeluaranFile[listing options={style=ITERAPlainKecil}]{keluaran/P12/k002.txt}
\begin{Catatan}
Batas dalam \texttt{n - 1 - i}: setiap putaran, satu elemen di ujung sudah final. \texttt{j + 1} tidak boleh melewati n − 1.
\end{Catatan}
}
\end{column}
\end{columns}
\note<1>{Minta mahasiswa menebak keluaran sebelum membuka slide berikutnya. Tanyakan satu atau dua tebakan yang berbeda, lalu buktikan.}
\end{frame}

\begin{frame}[fragile]{Selection sort}
\framesubtitle{Pilih yang terkecil}
\begin{columns}[T]
\begin{column}{0.60\textwidth}
\begin{Kode}[listing options={style=ITERACodeKecil}]{selection.cpp}
int a[5] = {29, 10, 14, 37, 13};
for (int i = 0; i < 4; i++) {
    int iMin = i;
    for (int j = i + 1; j < 5; j++) {
        if (a[j] < a[iMin]) iMin = j;
    }
    int t = a[i]; a[i] = a[iMin]; a[iMin] = t;
}
for (int k = 0; k < 5; k++) cout << a[k] << " ";
cout << endl;
\end{Kode}
\end{column}
\begin{column}{0.36\textwidth}
\only<1>{\begin{TebakDulu}
{\ITERAKickerFont\fontsize{10pt}{12pt}\selectfont\color{ITERAAlert}TEBAK DULU}\par\vspace{4pt}
Sebelum dijalankan, tulis keluarannya di kertas.
\end{TebakDulu}
}\only<2>{\KeluaranFile[listing options={style=ITERAPlainKecil}]{keluaran/P12/k003.txt}
\begin{Catatan}
Paling banyak satu pertukaran per putaran. Untuk urutan turun, cari yang terbesar.
\end{Catatan}
}
\end{column}
\end{columns}
\note<1>{Minta mahasiswa menebak keluaran sebelum membuka slide berikutnya. Tanyakan satu atau dua tebakan yang berbeda, lalu buktikan.}
\end{frame}

\begin{frame}[fragile]{Mengurutkan dua array sejajar}
\framesubtitle{Nama ikut skornya}
\begin{columns}[T]
\begin{column}{0.60\textwidth}
\begin{Kode}[listing options={style=ITERACodeKecil}]{sejajar.cpp}
string nama[4] = {"Rani", "Budi", "Siti", "Andi"};
int skor[4] = {82, 75, 95, 88};
for (int i = 0; i < 3; i++) {
    int iMaks = i;
    for (int j = i + 1; j < 4; j++)
        if (skor[j] > skor[iMaks]) iMaks = j;
    swap(skor[i], skor[iMaks]);
    swap(nama[i], nama[iMaks]);
}
for (int i = 0; i < 4; i++)
    cout << nama[i] << " " << skor[i] << endl;
\end{Kode}
\end{column}
\begin{column}{0.36\textwidth}
\only<1>{\begin{TebakDulu}
{\ITERAKickerFont\fontsize{10pt}{12pt}\selectfont\color{ITERAAlert}TEBAK DULU}\par\vspace{4pt}
Sebelum dijalankan, tulis keluarannya di kertas.
\end{TebakDulu}
}\only<2>{\KeluaranFile[listing options={style=ITERAPlainKecil}]{keluaran/P12/k004.txt}
\begin{Catatan}
Setiap pertukaran pada array kunci harus diikuti pertukaran yang sama pada array pasangannya. \texttt{swap} dari pustaka standar menukar dua nilai.
\end{Catatan}
}
\end{column}
\end{columns}
\note<1>{Minta mahasiswa menebak keluaran sebelum membuka slide berikutnya. Tanyakan satu atau dua tebakan yang berbeda, lalu buktikan.}
\end{frame}

\begin{frame}[fragile]{Tebak keluarannya}
\framesubtitle{Latihan menjelaskan, 3 menit}
\begin{columns}[T]
\begin{column}{0.56\textwidth}
\begin{Kode}[listing options={style=ITERACodeKecil}]{tebak.cpp}
int a[4] = {4, 3, 2, 1};
int tukar = 0;
for (int i = 0; i < 1; i++) {
    for (int j = 0; j < 3 - i; j++) {
        if (a[j] > a[j + 1]) {
            swap(a[j], a[j + 1]);
            tukar++;
        }
    }
}
for (int k = 0; k < 4; k++) cout << a[k];
cout << " " << tukar << endl;
\end{Kode}
\end{column}
\begin{column}{0.40\textwidth}
\begin{Catatan}
Tulis keluarannya di kertas, karakter demi karakter. Jangan dijalankan dulu.
\end{Catatan}
\end{column}
\end{columns}
\note{Contoh soal Bagian B. Beri waktu 3 menit sebelum slide jawaban.}
\end{frame}

\begin{frame}[fragile]{Tebak keluarannya: jawaban}
\framesubtitle{Latihan menjelaskan}
\begin{columns}[T]
\begin{column}{0.56\textwidth}
\begin{Kode}[listing options={style=ITERACodeKecil}]{tebak.cpp}
int a[4] = {4, 3, 2, 1};
int tukar = 0;
for (int i = 0; i < 1; i++) {
    for (int j = 0; j < 3 - i; j++) {
        if (a[j] > a[j + 1]) {
            swap(a[j], a[j + 1]);
            tukar++;
        }
    }
}
for (int k = 0; k < 4; k++) cout << a[k];
cout << " " << tukar << endl;
\end{Kode}
\end{column}
\begin{column}{0.40\textwidth}
\begin{Keluaran}[]
3214 3
\end{Keluaran}
\begin{Catatan}
Hanya satu putaran (i = 0). Angka 4 terdorong ke ujung lewat tiga pertukaran, tetapi sisanya belum terurut: 3 2 1 4.
\end{Catatan}
\end{column}
\end{columns}
\note{Tanyakan jawaban yang paling sering salah, lalu telusuri bersama.}
\end{frame}

\begin{frame}{Error yang sering muncul minggu ini}
\framesubtitle{Pesan diambil langsung dari g++}
\footnotesize
\begin{tabular}{@{}>{\raggedright\arraybackslash}p{208pt}>{\raggedright\arraybackslash}p{256pt}>{\raggedright\arraybackslash}p{398pt}@{}}
\kepala{Kesalahan} & \kepala{Contoh} & \kepala{Pesan pertama dari g++}\\
swap tanpa nilai yang bisa diubah & \texttt{swap(a[0], 5)} & \texttt{error: no matching function for call to 'swap(int\&, int)'}\\
\barisgenap Tukar string dengan int & \texttt{swap(nama[i], skor[i])} & \texttt{error: no matching function for call to 'swap(std::string\&, int\&)'}\\
Variabel sementara tak dipakai & \texttt{int t; tidak dipakai} & \texttt{warning: unused variable 't'}\\
\barisgenap Indeks bukan bilangan bulat & \texttt{a[j + 0.5]} & \texttt{error: invalid types 'int [3][double]' for array subscript}\\
Ukuran array tidak cocok & \texttt{int a[3] = \{1, 2, 3, 4\}} & \texttt{error: too many initializers for 'int [3]'}\\
\garisdasar
\end{tabular}
\vspace{10pt}

{\small Pesan di tabel ini dihasilkan dengan g++ dan opsi -Wall, sama seperti di cek.sh.}
\note{Minta mahasiswa memotret slide ini.}
\end{frame}

\Bagian{4}{Hands-on bertingkat}{45 menit, dari dasar sampai tantangan}
\note{Mahasiswa membuka folder 02 Hands-on/P12 - Sorting - Hands-on.}

\begin{frame}{Peta hands-on}
\framesubtitle{Folder 02 Hands-on/P12 - Sorting - Hands-on}
\small
\begin{tabular}{@{}>{\raggedright\arraybackslash}p{36pt}>{\raggedright\arraybackslash}p{267pt}>{\raggedright\arraybackslash}p{110pt}>{\raggedright\arraybackslash}p{83pt}>{\raggedright\arraybackslash}p{341pt}@{}}
\kepala{No} & \kepala{File} & \kepala{Tingkat} & \kepala{Waktu} & \kepala{Yang dilatih}\\
1 & \texttt{handson1-bubble.cpp} & Dasar & 10' & bubble sort naik dengan jejak per putaran\\
\barisgenap 2 & \texttt{handson2-peringkat.cpp} & Menengah & 15' & selection sort turun, dua array sejajar\\
3 & \texttt{handson3-tukar.cpp} & Tantangan & 20' & telusuri pertukaran, perbaiki tiga kesalahan\\
\barisgenap + & \texttt{bonus-hitung.cpp} & Pengayaan & sisa & bubble sort berhenti dini, hitung kerja\\
\garisdasar
\end{tabular}
\vspace{10pt}

Kerjakan berurutan. Cocokkan keluaran dengan folder \texttt{expected/}; latihan yang membaca masukan memakai data di folder \texttt{input/}.\note{Saat berkeliling, minta mahasiswa yang mengerjakan latihan tantangan menjelaskan blok PENJELASAN mereka.}
\end{frame}

\begin{frame}{Cara mengecek dan aturannya}
\framesubtitle{Hands-on}
\begin{columns}[T]
\begin{column}{0.47\textwidth}
\begin{Blok}{Di GDB Online}
Salin isi file latihan, lengkapi TODO, klik Run. Bila program membaca masukan, ketik isi file di \texttt{input/}. Bandingkan dengan \texttt{expected/}.
\end{Blok}
\end{column}
\begin{column}{0.47\textwidth}
\begin{BlokGelap}{Di komputer sendiri}
Jalankan \texttt{bash cek.sh}. Skrip mengompilasi dengan \texttt{-Wall -Werror}, memberi masukan dari \texttt{input/}, lalu mencocokkan keluaran.
\end{BlokGelap}
\end{column}
\end{columns}
\vspace{6pt}
\begin{Catatan}
AI boleh untuk memahami pesan error, tidak untuk meminta solusi utuh. Exit Ticket dikerjakan tanpa bantuan apa pun.
\end{Catatan}
\note{Hands-on tidak dinilai langsung; yang dinilai Exit Ticket dan tugas.}
\end{frame}

\Bagian{5}{Exit Ticket dan tugas}{25 menit terakhir, lalu pekerjaan rumah}
\note{Mulai Exit Ticket tepat waktu.}

\begin{frame}{Exit Ticket 12}
\framesubtitle{Individual, 25 menit, tanpa AI dan internet}
\small
\begin{tabular}{@{}>{\raggedright\arraybackslash}p{60pt}>{\raggedright\arraybackslash}p{560pt}>{\raggedright\arraybackslash}p{80pt}>{\raggedright\arraybackslash}p{150pt}@{}}
\kepala{Soal} & \kepala{Isi} & \kepala{Poin} & \kepala{CPMK}\\
A1 & Tulis bubble sort yang mengurutkan array turun & 25 & CPMK0607\\
\barisgenap A2 & Urutkan dua array sejajar (nama dan nilai) berdasarkan nilai & 25 & CPMK0607\\
B1 & Tuliskan isi array sesudah setiap putaran selection sort & 20 & CPMK0608\\
\barisgenap B2 & Jelaskan kesalahan pada potongan kode pertukaran & 20 & CPMK0608\\
B3 & Bandingkan banyak pertukaran bubble dan selection sort untuk data yang diberikan & 10 & CPMK0608\\
\garisdasar
\end{tabular}
\vspace{10pt}

{\small Soal A dinilai dari keluaran yang tepat sama. Soal B dinilai dari ketepatan dan alasan. Pelanggaran aturan membuat nilai Exit Ticket ini 0.}
\note{Soal lengkap dibagikan terpisah dan tidak disimpan di repo publik.}
\end{frame}

\begin{frame}{Tugas 12: Sistem Leaderboard}
\framesubtitle{Dikumpulkan sebelum Pertemuan 13}
\begin{columns}[T]
\begin{column}{0.47\textwidth}
\begin{Blok}{Bagian A, program, 50 poin}
Buat program leaderboard dengan menu berikut untuk maksimal 10 data (nama satu kata dan nilai). Ketentuannya di slide berikut.
\end{Blok}
\end{column}
\begin{column}{0.47\textwidth}
\begin{BlokGelap}{Bagian B, penjelasan, 50 poin}
Komentar maksimal 150 kata dan tabel isi array sesudah setiap putaran selection sort untuk lima data contoh. Jelaskan mengapa nama harus ikut ditukar, dan mengapa binary search hanya dipakai sesudah pengurutan abjad.
\end{BlokGelap}
\end{column}
\end{columns}
\vspace{8pt}
{\small Data di tugas adalah data kalian sendiri. Rubrik lengkap ada di Modul Belajar 12.}
\note{Ingatkan bahwa penjelasan disalin dari AI mudah dikenali karena tidak menyebut pengalaman mereka sendiri.}
\end{frame}

\begin{frame}{Ketentuan Tugas 12}
\framesubtitle{Sistem Leaderboard}
\footnotesize
\begin{tabular}{@{}>{\raggedright\arraybackslash}p{87pt}>{\raggedright\arraybackslash}p{788pt}@{}}
\kepala{No} & \kepala{Ketentuan Bagian A}\\
1 & Tambah data peserta dengan validasi nilai 0 sampai 100\\
\barisgenap 2 & Urutkan berdasarkan nilai turun dengan selection sort, lalu tampilkan leaderboard\\
3 & Urutkan berdasarkan nama (abjad) dengan bubble sort\\
\barisgenap 4 & Cari nama dengan linear search; tampilkan peringkat dan nilainya\\
5 & Cari nama dengan binary search pada data yang sudah urut abjad\\
\barisgenap 6 & Tampilkan banyak perbandingan dan pertukaran setiap algoritma yang dipakai, lalu keluar\\
Bonus & Bandingkan banyak langkah linear dan binary search untuk data yang sama, dan tampilkan dalam satu tabel.\\
\garisdasar
\end{tabular}
\note{Bacakan ketentuan satu per satu dan tanyakan apakah ada yang belum jelas.}
\end{frame}

\begin{frame}{Rubrik Tugas 12}
\framesubtitle{Empat level, sama di semua instrumen}
\footnotesize
\begin{tabular}{@{}>{\raggedright\arraybackslash}p{166pt}>{\raggedright\arraybackslash}p{44pt}>{\raggedright\arraybackslash}p{166pt}>{\raggedright\arraybackslash}p{156pt}>{\raggedright\arraybackslash}p{146pt}>{\raggedright\arraybackslash}p{146pt}@{}}
\kepala{Kriteria} & \kepala{Poin} & \kepala{Sangat baik 85--100} & \kepala{Baik 70--84} & \kepala{Cukup 55--69} & \kepala{Kurang <55}\\
A1 Program berjalan & 20 & tanpa error dan peringatan & ada peringatan & perlu satu perbaikan & tidak terkompilasi\\
\barisgenap A2 Kelengkapan menu & 15 & dua sort, dua search, penghitung & satu fitur kurang & dua kurang & tiga atau lebih kurang\\
A3 Ketepatan urutan dan pencarian & 15 & semua tepat termasuk nilai sama & satu kasus salah & dua salah & sebagian besar salah\\
\barisgenap B1 Jejak putaran dan alasan & 30 & tabel tepat dan alasan jelas & satu putaran keliru & tanpa alasan & tidak ada atau disalin\\
B2 Cerita error dan pengujian & 20 & pesan, sebab, perbaikan, uji & pesan tidak dikutip & hanya perbaikan & tidak ada\\
\garisdasar
\end{tabular}
\vspace{8pt}

{\small Nilai kriteria = poin × skor level. Bagian A total 50, Bagian B total 50.}
\note{Rubrik ini sama strukturnya setiap minggu; yang berubah hanya isi kriteria A2, A3, dan B1.}
\end{frame}

\begin{frame}{Sebelum Pertemuan 13}
\framesubtitle{Persiapan}
\begin{columns}[T]
\begin{column}{0.31\textwidth}
\begin{Langkah}{1}{Selesaikan}
Hands-on yang belum lulus \texttt{cek.sh}, terutama latihan tantangan.
\end{Langkah}
\end{column}
\begin{column}{0.31\textwidth}
\begin{Langkah}{2}{Kumpulkan}
Tugas 12 sebelum Pertemuan 13 dimulai.
\end{Langkah}
\end{column}
\begin{column}{0.31\textwidth}
\begin{Langkah}{3}{Baca}
Modul Belajar 12 untuk mengulang, lalu bagian awal modul berikutnya.
\end{Langkah}
\end{column}
\end{columns}
\vspace{10pt}
Pertemuan 13 membahas fungsi: memecah program menjadi bagian bernama yang bisa dipanggil berulang, termasuk fungsi pengurutan dan pencarian.\note{Tanyakan satu hal yang paling membingungkan hari ini untuk dibuka di review minggu depan.}
\end{frame}

\SlidePenutup{Terima kasih}{Sampai jumpa di Pertemuan 13: fungsi dan modularitas}
\note{Slide penutup.}
\end{document}
"
\documentclass[t]{beamer}
\usetheme{ITERA}
% Mode catatan pembicara: aktif bila \catatan didefinisikan sebelum \input.
\ifdefined\catatan
  \setbeameroption{show only notes}
  \setbeamertemplate{note page}{\vspace*{20pt}\hspace*{4pt}\begin{minipage}{900pt}
    {\ITERAKickerFont\fontsize{11pt}{14pt}\selectfont\color{ITERAGoldText}CATATAN PEMBICARA \ \ |\ \ SLIDE \insertframenumber}\par\vspace{4pt}
    {\ITERADisplay\bfseries\fontsize{22pt}{28pt}\selectfont\color{ITERAStructure}\insertframetitle}\par\vspace{12pt}
    \fontsize{15pt}{23pt}\selectfont\insertnote\end{minipage}}
\fi
\tikzset{fase/.style={draw=none,fill=#1,rounded corners=3pt,minimum width=158pt,minimum height=118pt,text width=146pt,align=center}}

\renewcommand\ITERAmeeting{Pertemuan 12}
\begin{document}
\SlideJudul{IF25-11002 PRAKTIKUM PEMROGRAMAN}{Pertemuan 12\\Sorting}{Mengurutkan data dengan bubble sort dan selection sort, serta menelusuri setiap putarannya}{Tim Pengajar}{Tahun 2026. Sejajar dengan Pertemuan 12 Algoritma Pemrograman: Sorting.}
\note{Buka dengan menanyakan satu hal dari pertemuan lalu yang masih membingungkan.}

\begin{frame}{Dua mata kuliah, satu minggu}
\framesubtitle{Sebelum mulai}
Topik minggu ini sama di dua mata kuliah, dengan peran yang berbeda.
\vspace{10pt}

\begin{columns}[T]
\begin{column}{0.47\textwidth}
\begin{Blok}{Algoritma: Sorting}
Algoritma pengurutan gelembung dan seleksi, menelusuri isi larik per putaran, dan menghitung perbandingan serta pertukaran.
\end{Blok}
\end{column}
\begin{column}{0.47\textwidth}
\begin{BlokGelap}{Praktikum: Sorting}
Menulis pertukaran, bubble sort, dan selection sort di C++, termasuk mengurutkan dua array sejajar.
\end{BlokGelap}
\end{column}
\end{columns}
\note{Tanyakan apa yang sudah dibahas di Algoritma minggu ini. Gunakan jawabannya sebagai jembatan ke praktikum.}
\end{frame}

\begin{frame}{Saat pulang nanti, kalian bisa}
\framesubtitle{Target hari ini}
\begin{columns}[T]
\begin{column}{0.47\textwidth}
\begin{Langkah}{1}{Menerapkan}
Menulis program yang mengurutkan array naik maupun turun dengan bubble sort dan selection sort, termasuk mengurutkan dua array sejajar sekaligus

{\footnotesize\color{ITERAInkMuted}Sub-CPMK 12.1, CPMK0607}
\end{Langkah}
\end{column}
\begin{column}{0.47\textwidth}
\begin{Langkah}{2}{Menjelaskan}
Menelusuri isi array sesudah setiap putaran pengurutan dan menjelaskan banyak perbandingan serta pertukaran yang terjadi

{\footnotesize\color{ITERAInkMuted}Sub-CPMK 12.2, CPMK0608}
\end{Langkah}
\end{column}
\end{columns}
\vspace{8pt}
\begin{Catatan}
Dua kemampuan ini dinilai sama berat di Exit Ticket, tugas, UTS, dan UAS.
\end{Catatan}
\note{Bacakan target dengan pelan. Kata kuncinya menerapkan dan menjelaskan.}
\end{frame}

\begin{frame}{Ingat minggu lalu}
\framesubtitle{Review, 5 menit}
\begin{columns}[T]
\begin{column}{0.31\textwidth}
\begin{BlokGelap}{Pertanyaan 1}
Apa prasyarat binary search?
\end{BlokGelap}
\end{column}
\begin{column}{0.31\textwidth}
\begin{BlokGelap}{Pertanyaan 2}
Mengapa syarat binary search \texttt{low <= high}?
\end{BlokGelap}
\end{column}
\begin{column}{0.31\textwidth}
\begin{BlokGelap}{Pertanyaan 3}
Berapa langkah binary search untuk 1.000 data?
\end{BlokGelap}
\end{column}
\end{columns}
\vspace{8pt}
Jawab di kertas dulu, lalu cocokkan dengan teman sebelah.\note{Pilih dua mahasiswa untuk menjawab. Koreksi singkat, jangan mengajar ulang.}
\end{frame}

\Bagian{1}{Masalah pembuka}{Papan peringkat lomba coding}
\note{Masalah pembuka 10 menit.}

\begin{frame}[fragile]{Papan peringkat lomba coding}
\framesubtitle{Masalah minggu ini}
\begin{columns}[T]
\begin{column}{0.55\textwidth}
{\small Panitia lomba coding mencatat skor peserta sesuai urutan pendaftaran. Di akhir lomba, papan peringkat harus menampilkan peserta dari skor tertinggi.

\vspace{6pt}
Mencari skor tertinggi sudah bisa sejak Pertemuan 9. Tantangannya sekarang adalah menyusun semua peserta, dan nama harus tetap menempel pada skornya.\par}
\vspace{6pt}
\begin{Catatan}
\textbf{Diskusi 2 menit.} Bagaimana kalian mengurutkan lima kartu skor di atas meja dengan tangan? Langkah apa yang terus kalian ulang?
\end{Catatan}
\end{column}
\begin{column}{0.41\textwidth}
\begin{Keluaran}[]
1. Siti 95
2. Andi 88
3. Rani 82
4. Budi 75
5. Doni 70
\end{Keluaran}
\end{column}
\end{columns}
\note{Tampung beberapa jawaban tanpa menilai benar salah.}
\end{frame}

\begin{frame}{Dari langkah ke instruksi}
\framesubtitle{Sambungan dengan Algoritma}
\begin{columns}[T]
\begin{column}{0.31\textwidth}
\begin{Langkah}{1}{Pilih caranya}
Tukar pasangan bersebelahan yang terbalik, atau cari yang terbesar lalu taruh di depan.
\end{Langkah}
\end{column}
\begin{column}{0.31\textwidth}
\begin{Langkah}{2}{Ulangi per putaran}
Setiap putaran memastikan satu elemen sudah di tempat akhirnya.
\end{Langkah}
\end{column}
\begin{column}{0.31\textwidth}
\begin{Langkah}{3}{Terjemahkan}
Pertukaran tiga baris dan dua \texttt{for} bersarang.
\end{Langkah}
\end{column}
\end{columns}
\vspace{10pt}
Langkah 1 dan 2 dilatih di Algoritma. Langkah 3 kita kerjakan hari ini dengan C++.\note{Hubungkan eksplisit ke Algoritma minggu ini.}
\end{frame}

\Bagian{2}{Coba dulu, baru dijelaskan}{25 menit eksperimen di GDB Online}
\note{Struggle 25 menit. Berkeliling, jangan menjelaskan materi dulu.}

\begin{frame}{Misi kalian}
\framesubtitle{Struggle, 25 menit}
\begin{columns}[T]
\begin{column}{0.31\textwidth}
\begin{Langkah}{1}{Tukar}
Tukar isi \texttt{a} dan \texttt{b} sehingga \texttt{a = 7, b = 3} menjadi \texttt{a = 3, b = 7}.
\end{Langkah}
\end{column}
\begin{column}{0.31\textwidth}
\begin{Langkah}{2}{Satu putaran}
Pada \texttt{\{5, 1, 4, 2\}}, tukar setiap pasangan bersebelahan yang terbalik dari kiri ke kanan.
\end{Langkah}
\end{column}
\begin{column}{0.31\textwidth}
\begin{Langkah}{3}{Semua putaran}
Ulangi putaran itu sampai data terurut. Berapa putaran yang dibutuhkan?
\end{Langkah}
\end{column}
\end{columns}
\vspace{8pt}
\begin{columns}[T]
\begin{column}{0.47\textwidth}
\begin{Blok}{Boleh}
Bertanya ke teman, mengubah apa saja, sengaja membuat error.
\end{Blok}
\end{column}
\begin{column}{0.47\textwidth}
\begin{BlokAlert}{Belum dulu}
Mencari jawaban jadi di internet atau AI.
\end{BlokAlert}
\end{column}
\end{columns}
\note{Catat kesalahan yang sering muncul untuk dibahas di debrief.}
\end{frame}

\begin{frame}{Apa yang kalian temukan?}
\framesubtitle{Debrief struggle}
\begin{columns}[T]
\begin{column}{0.31\textwidth}
\begin{BlokGelap}{Pertanyaan 1}
Apa yang terjadi bila menukar dengan \texttt{a = b; b = a;}?
\end{BlokGelap}
\end{column}
\begin{column}{0.31\textwidth}
\begin{BlokGelap}{Pertanyaan 2}
Sesudah satu putaran, elemen mana yang pasti sudah di tempatnya?
\end{BlokGelap}
\end{column}
\begin{column}{0.31\textwidth}
\begin{BlokGelap}{Pertanyaan 3}
Apakah selalu butuh n − 1 putaran?
\end{BlokGelap}
\end{column}
\end{columns}
\vspace{10pt}
Jawaban kalian akan kita uji di bagian berikutnya.\note{Tulis jawaban di papan; buktikan atau patahkan di bagian materi.}
\end{frame}

\Bagian{3}{Mengurutkan langkah demi langkah}{Pertukaran, bubble sort, selection sort, dan array sejajar}
\note{Materi 40 menit. Tunjukkan setiap contoh langsung di GDB Online.}

\begin{frame}[fragile]{Menukar dua nilai}
\framesubtitle{Variabel sementara}
\begin{columns}[T]
\begin{column}{0.55\textwidth}
\begin{Kode}[]{tukar.cpp}
int a = 7, b = 3;
int t = a;
a = b;
b = t;
cout << a << " " << b << endl;
\end{Kode}
\end{column}
\begin{column}{0.41\textwidth}
\only<1>{\begin{TebakDulu}
{\ITERAKickerFont\fontsize{10pt}{12pt}\selectfont\color{ITERAAlert}TEBAK DULU}\par\vspace{4pt}
Sebelum dijalankan, tulis keluarannya di kertas.
\end{TebakDulu}
}\only<2>{\KeluaranFile[]{keluaran/P12/k001.txt}
\begin{Catatan}
Tanpa \texttt{t}, nilai lama \texttt{a} hilang begitu ditimpa.
\end{Catatan}
}
\end{column}
\end{columns}
\note<1>{Minta mahasiswa menebak keluaran sebelum membuka slide berikutnya. Tanyakan satu atau dua tebakan yang berbeda, lalu buktikan.}
\end{frame}

\begin{frame}{Bubble sort: satu putaran}
\framesubtitle{Pasangan bersebelahan}
Putaran pertama pada \{5, 1, 4, 2, 8\}: bandingkan setiap pasangan, tukar bila terbalik.
\vspace{8pt}

\small
\begin{tabular}{@{}>{\raggedright\arraybackslash}p{172pt}>{\raggedright\arraybackslash}p{326pt}>{\raggedright\arraybackslash}p{364pt}@{}}
\kepala{Pasangan} & \kepala{Isi array sesudahnya} & \kepala{Keterangan}\\
a[0], a[1] & 1 5 4 2 8 & 5 > 1, tukar\\
\barisgenap a[1], a[2] & 1 4 5 2 8 & 5 > 4, tukar\\
a[2], a[3] & 1 4 2 5 8 & 5 > 2, tukar\\
\barisgenap a[3], a[4] & 1 4 2 5 8 & 5 < 8, tetap\\
\garisdasar
\end{tabular}
\vspace{10pt}

Sesudah putaran pertama, elemen terbesar pasti sudah di ujung kanan, seperti gelembung yang naik ke permukaan.
\end{frame}

\begin{frame}[fragile]{Bubble sort di C++}
\framesubtitle{Dua perulangan bersarang}
\begin{columns}[T]
\begin{column}{0.60\textwidth}
\begin{Kode}[listing options={style=ITERACodeKecil}]{bubble.cpp}
int a[5] = {5, 1, 4, 2, 8};
int n = 5;
for (int i = 0; i < n - 1; i++) {
    for (int j = 0; j < n - 1 - i; j++) {
        if (a[j] > a[j + 1]) {
            int t = a[j]; a[j] = a[j + 1]; a[j + 1] = t;
        }
    }
}
for (int k = 0; k < n; k++) cout << a[k] << " ";
cout << endl;
\end{Kode}
\end{column}
\begin{column}{0.36\textwidth}
\only<1>{\begin{TebakDulu}
{\ITERAKickerFont\fontsize{10pt}{12pt}\selectfont\color{ITERAAlert}TEBAK DULU}\par\vspace{4pt}
Sebelum dijalankan, tulis keluarannya di kertas.
\end{TebakDulu}
}\only<2>{\KeluaranFile[listing options={style=ITERAPlainKecil}]{keluaran/P12/k002.txt}
\begin{Catatan}
Batas dalam \texttt{n - 1 - i}: setiap putaran, satu elemen di ujung sudah final. \texttt{j + 1} tidak boleh melewati n − 1.
\end{Catatan}
}
\end{column}
\end{columns}
\note<1>{Minta mahasiswa menebak keluaran sebelum membuka slide berikutnya. Tanyakan satu atau dua tebakan yang berbeda, lalu buktikan.}
\end{frame}

\begin{frame}[fragile]{Selection sort}
\framesubtitle{Pilih yang terkecil}
\begin{columns}[T]
\begin{column}{0.60\textwidth}
\begin{Kode}[listing options={style=ITERACodeKecil}]{selection.cpp}
int a[5] = {29, 10, 14, 37, 13};
for (int i = 0; i < 4; i++) {
    int iMin = i;
    for (int j = i + 1; j < 5; j++) {
        if (a[j] < a[iMin]) iMin = j;
    }
    int t = a[i]; a[i] = a[iMin]; a[iMin] = t;
}
for (int k = 0; k < 5; k++) cout << a[k] << " ";
cout << endl;
\end{Kode}
\end{column}
\begin{column}{0.36\textwidth}
\only<1>{\begin{TebakDulu}
{\ITERAKickerFont\fontsize{10pt}{12pt}\selectfont\color{ITERAAlert}TEBAK DULU}\par\vspace{4pt}
Sebelum dijalankan, tulis keluarannya di kertas.
\end{TebakDulu}
}\only<2>{\KeluaranFile[listing options={style=ITERAPlainKecil}]{keluaran/P12/k003.txt}
\begin{Catatan}
Paling banyak satu pertukaran per putaran. Untuk urutan turun, cari yang terbesar.
\end{Catatan}
}
\end{column}
\end{columns}
\note<1>{Minta mahasiswa menebak keluaran sebelum membuka slide berikutnya. Tanyakan satu atau dua tebakan yang berbeda, lalu buktikan.}
\end{frame}

\begin{frame}[fragile]{Mengurutkan dua array sejajar}
\framesubtitle{Nama ikut skornya}
\begin{columns}[T]
\begin{column}{0.60\textwidth}
\begin{Kode}[listing options={style=ITERACodeKecil}]{sejajar.cpp}
string nama[4] = {"Rani", "Budi", "Siti", "Andi"};
int skor[4] = {82, 75, 95, 88};
for (int i = 0; i < 3; i++) {
    int iMaks = i;
    for (int j = i + 1; j < 4; j++)
        if (skor[j] > skor[iMaks]) iMaks = j;
    swap(skor[i], skor[iMaks]);
    swap(nama[i], nama[iMaks]);
}
for (int i = 0; i < 4; i++)
    cout << nama[i] << " " << skor[i] << endl;
\end{Kode}
\end{column}
\begin{column}{0.36\textwidth}
\only<1>{\begin{TebakDulu}
{\ITERAKickerFont\fontsize{10pt}{12pt}\selectfont\color{ITERAAlert}TEBAK DULU}\par\vspace{4pt}
Sebelum dijalankan, tulis keluarannya di kertas.
\end{TebakDulu}
}\only<2>{\KeluaranFile[listing options={style=ITERAPlainKecil}]{keluaran/P12/k004.txt}
\begin{Catatan}
Setiap pertukaran pada array kunci harus diikuti pertukaran yang sama pada array pasangannya. \texttt{swap} dari pustaka standar menukar dua nilai.
\end{Catatan}
}
\end{column}
\end{columns}
\note<1>{Minta mahasiswa menebak keluaran sebelum membuka slide berikutnya. Tanyakan satu atau dua tebakan yang berbeda, lalu buktikan.}
\end{frame}

\begin{frame}[fragile]{Tebak keluarannya}
\framesubtitle{Latihan menjelaskan, 3 menit}
\begin{columns}[T]
\begin{column}{0.56\textwidth}
\begin{Kode}[listing options={style=ITERACodeKecil}]{tebak.cpp}
int a[4] = {4, 3, 2, 1};
int tukar = 0;
for (int i = 0; i < 1; i++) {
    for (int j = 0; j < 3 - i; j++) {
        if (a[j] > a[j + 1]) {
            swap(a[j], a[j + 1]);
            tukar++;
        }
    }
}
for (int k = 0; k < 4; k++) cout << a[k];
cout << " " << tukar << endl;
\end{Kode}
\end{column}
\begin{column}{0.40\textwidth}
\begin{Catatan}
Tulis keluarannya di kertas, karakter demi karakter. Jangan dijalankan dulu.
\end{Catatan}
\end{column}
\end{columns}
\note{Contoh soal Bagian B. Beri waktu 3 menit sebelum slide jawaban.}
\end{frame}

\begin{frame}[fragile]{Tebak keluarannya: jawaban}
\framesubtitle{Latihan menjelaskan}
\begin{columns}[T]
\begin{column}{0.56\textwidth}
\begin{Kode}[listing options={style=ITERACodeKecil}]{tebak.cpp}
int a[4] = {4, 3, 2, 1};
int tukar = 0;
for (int i = 0; i < 1; i++) {
    for (int j = 0; j < 3 - i; j++) {
        if (a[j] > a[j + 1]) {
            swap(a[j], a[j + 1]);
            tukar++;
        }
    }
}
for (int k = 0; k < 4; k++) cout << a[k];
cout << " " << tukar << endl;
\end{Kode}
\end{column}
\begin{column}{0.40\textwidth}
\begin{Keluaran}[]
3214 3
\end{Keluaran}
\begin{Catatan}
Hanya satu putaran (i = 0). Angka 4 terdorong ke ujung lewat tiga pertukaran, tetapi sisanya belum terurut: 3 2 1 4.
\end{Catatan}
\end{column}
\end{columns}
\note{Tanyakan jawaban yang paling sering salah, lalu telusuri bersama.}
\end{frame}

\begin{frame}{Error yang sering muncul minggu ini}
\framesubtitle{Pesan diambil langsung dari g++}
\footnotesize
\begin{tabular}{@{}>{\raggedright\arraybackslash}p{208pt}>{\raggedright\arraybackslash}p{256pt}>{\raggedright\arraybackslash}p{398pt}@{}}
\kepala{Kesalahan} & \kepala{Contoh} & \kepala{Pesan pertama dari g++}\\
swap tanpa nilai yang bisa diubah & \texttt{swap(a[0], 5)} & \texttt{error: no matching function for call to 'swap(int\&, int)'}\\
\barisgenap Tukar string dengan int & \texttt{swap(nama[i], skor[i])} & \texttt{error: no matching function for call to 'swap(std::string\&, int\&)'}\\
Variabel sementara tak dipakai & \texttt{int t; tidak dipakai} & \texttt{warning: unused variable 't'}\\
\barisgenap Indeks bukan bilangan bulat & \texttt{a[j + 0.5]} & \texttt{error: invalid types 'int [3][double]' for array subscript}\\
Ukuran array tidak cocok & \texttt{int a[3] = \{1, 2, 3, 4\}} & \texttt{error: too many initializers for 'int [3]'}\\
\garisdasar
\end{tabular}
\vspace{10pt}

{\small Pesan di tabel ini dihasilkan dengan g++ dan opsi -Wall, sama seperti di cek.sh.}
\note{Minta mahasiswa memotret slide ini.}
\end{frame}

\Bagian{4}{Hands-on bertingkat}{45 menit, dari dasar sampai tantangan}
\note{Mahasiswa membuka folder 02 Hands-on/P12 - Sorting - Hands-on.}

\begin{frame}{Peta hands-on}
\framesubtitle{Folder 02 Hands-on/P12 - Sorting - Hands-on}
\small
\begin{tabular}{@{}>{\raggedright\arraybackslash}p{36pt}>{\raggedright\arraybackslash}p{267pt}>{\raggedright\arraybackslash}p{110pt}>{\raggedright\arraybackslash}p{83pt}>{\raggedright\arraybackslash}p{341pt}@{}}
\kepala{No} & \kepala{File} & \kepala{Tingkat} & \kepala{Waktu} & \kepala{Yang dilatih}\\
1 & \texttt{handson1-bubble.cpp} & Dasar & 10' & bubble sort naik dengan jejak per putaran\\
\barisgenap 2 & \texttt{handson2-peringkat.cpp} & Menengah & 15' & selection sort turun, dua array sejajar\\
3 & \texttt{handson3-tukar.cpp} & Tantangan & 20' & telusuri pertukaran, perbaiki tiga kesalahan\\
\barisgenap + & \texttt{bonus-hitung.cpp} & Pengayaan & sisa & bubble sort berhenti dini, hitung kerja\\
\garisdasar
\end{tabular}
\vspace{10pt}

Kerjakan berurutan. Cocokkan keluaran dengan folder \texttt{expected/}; latihan yang membaca masukan memakai data di folder \texttt{input/}.\note{Saat berkeliling, minta mahasiswa yang mengerjakan latihan tantangan menjelaskan blok PENJELASAN mereka.}
\end{frame}

\begin{frame}{Cara mengecek dan aturannya}
\framesubtitle{Hands-on}
\begin{columns}[T]
\begin{column}{0.47\textwidth}
\begin{Blok}{Di GDB Online}
Salin isi file latihan, lengkapi TODO, klik Run. Bila program membaca masukan, ketik isi file di \texttt{input/}. Bandingkan dengan \texttt{expected/}.
\end{Blok}
\end{column}
\begin{column}{0.47\textwidth}
\begin{BlokGelap}{Di komputer sendiri}
Jalankan \texttt{bash cek.sh}. Skrip mengompilasi dengan \texttt{-Wall -Werror}, memberi masukan dari \texttt{input/}, lalu mencocokkan keluaran.
\end{BlokGelap}
\end{column}
\end{columns}
\vspace{6pt}
\begin{Catatan}
AI boleh untuk memahami pesan error, tidak untuk meminta solusi utuh. Exit Ticket dikerjakan tanpa bantuan apa pun.
\end{Catatan}
\note{Hands-on tidak dinilai langsung; yang dinilai Exit Ticket dan tugas.}
\end{frame}

\Bagian{5}{Exit Ticket dan tugas}{25 menit terakhir, lalu pekerjaan rumah}
\note{Mulai Exit Ticket tepat waktu.}

\begin{frame}{Exit Ticket 12}
\framesubtitle{Individual, 25 menit, tanpa AI dan internet}
\small
\begin{tabular}{@{}>{\raggedright\arraybackslash}p{60pt}>{\raggedright\arraybackslash}p{560pt}>{\raggedright\arraybackslash}p{80pt}>{\raggedright\arraybackslash}p{150pt}@{}}
\kepala{Soal} & \kepala{Isi} & \kepala{Poin} & \kepala{CPMK}\\
A1 & Tulis bubble sort yang mengurutkan array turun & 25 & CPMK0607\\
\barisgenap A2 & Urutkan dua array sejajar (nama dan nilai) berdasarkan nilai & 25 & CPMK0607\\
B1 & Tuliskan isi array sesudah setiap putaran selection sort & 20 & CPMK0608\\
\barisgenap B2 & Jelaskan kesalahan pada potongan kode pertukaran & 20 & CPMK0608\\
B3 & Bandingkan banyak pertukaran bubble dan selection sort untuk data yang diberikan & 10 & CPMK0608\\
\garisdasar
\end{tabular}
\vspace{10pt}

{\small Soal A dinilai dari keluaran yang tepat sama. Soal B dinilai dari ketepatan dan alasan. Pelanggaran aturan membuat nilai Exit Ticket ini 0.}
\note{Soal lengkap dibagikan terpisah dan tidak disimpan di repo publik.}
\end{frame}

\begin{frame}{Tugas 12: Sistem Leaderboard}
\framesubtitle{Dikumpulkan sebelum Pertemuan 13}
\begin{columns}[T]
\begin{column}{0.47\textwidth}
\begin{Blok}{Bagian A, program, 50 poin}
Buat program leaderboard dengan menu berikut untuk maksimal 10 data (nama satu kata dan nilai). Ketentuannya di slide berikut.
\end{Blok}
\end{column}
\begin{column}{0.47\textwidth}
\begin{BlokGelap}{Bagian B, penjelasan, 50 poin}
Komentar maksimal 150 kata dan tabel isi array sesudah setiap putaran selection sort untuk lima data contoh. Jelaskan mengapa nama harus ikut ditukar, dan mengapa binary search hanya dipakai sesudah pengurutan abjad.
\end{BlokGelap}
\end{column}
\end{columns}
\vspace{8pt}
{\small Data di tugas adalah data kalian sendiri. Rubrik lengkap ada di Modul Belajar 12.}
\note{Ingatkan bahwa penjelasan disalin dari AI mudah dikenali karena tidak menyebut pengalaman mereka sendiri.}
\end{frame}

\begin{frame}{Ketentuan Tugas 12}
\framesubtitle{Sistem Leaderboard}
\footnotesize
\begin{tabular}{@{}>{\raggedright\arraybackslash}p{87pt}>{\raggedright\arraybackslash}p{788pt}@{}}
\kepala{No} & \kepala{Ketentuan Bagian A}\\
1 & Tambah data peserta dengan validasi nilai 0 sampai 100\\
\barisgenap 2 & Urutkan berdasarkan nilai turun dengan selection sort, lalu tampilkan leaderboard\\
3 & Urutkan berdasarkan nama (abjad) dengan bubble sort\\
\barisgenap 4 & Cari nama dengan linear search; tampilkan peringkat dan nilainya\\
5 & Cari nama dengan binary search pada data yang sudah urut abjad\\
\barisgenap 6 & Tampilkan banyak perbandingan dan pertukaran setiap algoritma yang dipakai, lalu keluar\\
Bonus & Bandingkan banyak langkah linear dan binary search untuk data yang sama, dan tampilkan dalam satu tabel.\\
\garisdasar
\end{tabular}
\note{Bacakan ketentuan satu per satu dan tanyakan apakah ada yang belum jelas.}
\end{frame}

\begin{frame}{Rubrik Tugas 12}
\framesubtitle{Empat level, sama di semua instrumen}
\footnotesize
\begin{tabular}{@{}>{\raggedright\arraybackslash}p{166pt}>{\raggedright\arraybackslash}p{44pt}>{\raggedright\arraybackslash}p{166pt}>{\raggedright\arraybackslash}p{156pt}>{\raggedright\arraybackslash}p{146pt}>{\raggedright\arraybackslash}p{146pt}@{}}
\kepala{Kriteria} & \kepala{Poin} & \kepala{Sangat baik 85--100} & \kepala{Baik 70--84} & \kepala{Cukup 55--69} & \kepala{Kurang <55}\\
A1 Program berjalan & 20 & tanpa error dan peringatan & ada peringatan & perlu satu perbaikan & tidak terkompilasi\\
\barisgenap A2 Kelengkapan menu & 15 & dua sort, dua search, penghitung & satu fitur kurang & dua kurang & tiga atau lebih kurang\\
A3 Ketepatan urutan dan pencarian & 15 & semua tepat termasuk nilai sama & satu kasus salah & dua salah & sebagian besar salah\\
\barisgenap B1 Jejak putaran dan alasan & 30 & tabel tepat dan alasan jelas & satu putaran keliru & tanpa alasan & tidak ada atau disalin\\
B2 Cerita error dan pengujian & 20 & pesan, sebab, perbaikan, uji & pesan tidak dikutip & hanya perbaikan & tidak ada\\
\garisdasar
\end{tabular}
\vspace{8pt}

{\small Nilai kriteria = poin × skor level. Bagian A total 50, Bagian B total 50.}
\note{Rubrik ini sama strukturnya setiap minggu; yang berubah hanya isi kriteria A2, A3, dan B1.}
\end{frame}

\begin{frame}{Sebelum Pertemuan 13}
\framesubtitle{Persiapan}
\begin{columns}[T]
\begin{column}{0.31\textwidth}
\begin{Langkah}{1}{Selesaikan}
Hands-on yang belum lulus \texttt{cek.sh}, terutama latihan tantangan.
\end{Langkah}
\end{column}
\begin{column}{0.31\textwidth}
\begin{Langkah}{2}{Kumpulkan}
Tugas 12 sebelum Pertemuan 13 dimulai.
\end{Langkah}
\end{column}
\begin{column}{0.31\textwidth}
\begin{Langkah}{3}{Baca}
Modul Belajar 12 untuk mengulang, lalu bagian awal modul berikutnya.
\end{Langkah}
\end{column}
\end{columns}
\vspace{10pt}
Pertemuan 13 membahas fungsi: memecah program menjadi bagian bernama yang bisa dipanggil berulang, termasuk fungsi pengurutan dan pencarian.\note{Tanyakan satu hal yang paling membingungkan hari ini untuk dibuka di review minggu depan.}
\end{frame}

\SlidePenutup{Terima kasih}{Sampai jumpa di Pertemuan 13: fungsi dan modularitas}
\note{Slide penutup.}
\end{document}
"
\documentclass[t]{beamer}
\usetheme{ITERA}
% Mode catatan pembicara: aktif bila \catatan didefinisikan sebelum \input.
\ifdefined\catatan
  \setbeameroption{show only notes}
  \setbeamertemplate{note page}{\vspace*{20pt}\hspace*{4pt}\begin{minipage}{900pt}
    {\ITERAKickerFont\fontsize{11pt}{14pt}\selectfont\color{ITERAGoldText}CATATAN PEMBICARA \ \ |\ \ SLIDE \insertframenumber}\par\vspace{4pt}
    {\ITERADisplay\bfseries\fontsize{22pt}{28pt}\selectfont\color{ITERAStructure}\insertframetitle}\par\vspace{12pt}
    \fontsize{15pt}{23pt}\selectfont\insertnote\end{minipage}}
\fi
\tikzset{fase/.style={draw=none,fill=#1,rounded corners=3pt,minimum width=158pt,minimum height=118pt,text width=146pt,align=center}}

\renewcommand\ITERAmeeting{Pertemuan 12}
\begin{document}
\SlideJudul{IF25-11002 PRAKTIKUM PEMROGRAMAN}{Pertemuan 12\\Sorting}{Mengurutkan data dengan bubble sort dan selection sort, serta menelusuri setiap putarannya}{Tim Pengajar}{Tahun 2026. Sejajar dengan Pertemuan 12 Algoritma Pemrograman: Sorting.}
\note{Buka dengan menanyakan satu hal dari pertemuan lalu yang masih membingungkan.}

\begin{frame}{Dua mata kuliah, satu minggu}
\framesubtitle{Sebelum mulai}
Topik minggu ini sama di dua mata kuliah, dengan peran yang berbeda.
\vspace{10pt}

\begin{columns}[T]
\begin{column}{0.47\textwidth}
\begin{Blok}{Algoritma: Sorting}
Algoritma pengurutan gelembung dan seleksi, menelusuri isi larik per putaran, dan menghitung perbandingan serta pertukaran.
\end{Blok}
\end{column}
\begin{column}{0.47\textwidth}
\begin{BlokGelap}{Praktikum: Sorting}
Menulis pertukaran, bubble sort, dan selection sort di C++, termasuk mengurutkan dua array sejajar.
\end{BlokGelap}
\end{column}
\end{columns}
\note{Tanyakan apa yang sudah dibahas di Algoritma minggu ini. Gunakan jawabannya sebagai jembatan ke praktikum.}
\end{frame}

\begin{frame}{Saat pulang nanti, kalian bisa}
\framesubtitle{Target hari ini}
\begin{columns}[T]
\begin{column}{0.47\textwidth}
\begin{Langkah}{1}{Menerapkan}
Menulis program yang mengurutkan array naik maupun turun dengan bubble sort dan selection sort, termasuk mengurutkan dua array sejajar sekaligus

{\footnotesize\color{ITERAInkMuted}Sub-CPMK 12.1, CPMK0607}
\end{Langkah}
\end{column}
\begin{column}{0.47\textwidth}
\begin{Langkah}{2}{Menjelaskan}
Menelusuri isi array sesudah setiap putaran pengurutan dan menjelaskan banyak perbandingan serta pertukaran yang terjadi

{\footnotesize\color{ITERAInkMuted}Sub-CPMK 12.2, CPMK0608}
\end{Langkah}
\end{column}
\end{columns}
\vspace{8pt}
\begin{Catatan}
Dua kemampuan ini dinilai sama berat di Exit Ticket, tugas, UTS, dan UAS.
\end{Catatan}
\note{Bacakan target dengan pelan. Kata kuncinya menerapkan dan menjelaskan.}
\end{frame}

\begin{frame}{Ingat minggu lalu}
\framesubtitle{Review, 5 menit}
\begin{columns}[T]
\begin{column}{0.31\textwidth}
\begin{BlokGelap}{Pertanyaan 1}
Apa prasyarat binary search?
\end{BlokGelap}
\end{column}
\begin{column}{0.31\textwidth}
\begin{BlokGelap}{Pertanyaan 2}
Mengapa syarat binary search \texttt{low <= high}?
\end{BlokGelap}
\end{column}
\begin{column}{0.31\textwidth}
\begin{BlokGelap}{Pertanyaan 3}
Berapa langkah binary search untuk 1.000 data?
\end{BlokGelap}
\end{column}
\end{columns}
\vspace{8pt}
Jawab di kertas dulu, lalu cocokkan dengan teman sebelah.\note{Pilih dua mahasiswa untuk menjawab. Koreksi singkat, jangan mengajar ulang.}
\end{frame}

\Bagian{1}{Masalah pembuka}{Papan peringkat lomba coding}
\note{Masalah pembuka 10 menit.}

\begin{frame}[fragile]{Papan peringkat lomba coding}
\framesubtitle{Masalah minggu ini}
\begin{columns}[T]
\begin{column}{0.55\textwidth}
{\small Panitia lomba coding mencatat skor peserta sesuai urutan pendaftaran. Di akhir lomba, papan peringkat harus menampilkan peserta dari skor tertinggi.

\vspace{6pt}
Mencari skor tertinggi sudah bisa sejak Pertemuan 9. Tantangannya sekarang adalah menyusun semua peserta, dan nama harus tetap menempel pada skornya.\par}
\vspace{6pt}
\begin{Catatan}
\textbf{Diskusi 2 menit.} Bagaimana kalian mengurutkan lima kartu skor di atas meja dengan tangan? Langkah apa yang terus kalian ulang?
\end{Catatan}
\end{column}
\begin{column}{0.41\textwidth}
\begin{Keluaran}[]
1. Siti 95
2. Andi 88
3. Rani 82
4. Budi 75
5. Doni 70
\end{Keluaran}
\end{column}
\end{columns}
\note{Tampung beberapa jawaban tanpa menilai benar salah.}
\end{frame}

\begin{frame}{Dari langkah ke instruksi}
\framesubtitle{Sambungan dengan Algoritma}
\begin{columns}[T]
\begin{column}{0.31\textwidth}
\begin{Langkah}{1}{Pilih caranya}
Tukar pasangan bersebelahan yang terbalik, atau cari yang terbesar lalu taruh di depan.
\end{Langkah}
\end{column}
\begin{column}{0.31\textwidth}
\begin{Langkah}{2}{Ulangi per putaran}
Setiap putaran memastikan satu elemen sudah di tempat akhirnya.
\end{Langkah}
\end{column}
\begin{column}{0.31\textwidth}
\begin{Langkah}{3}{Terjemahkan}
Pertukaran tiga baris dan dua \texttt{for} bersarang.
\end{Langkah}
\end{column}
\end{columns}
\vspace{10pt}
Langkah 1 dan 2 dilatih di Algoritma. Langkah 3 kita kerjakan hari ini dengan C++.\note{Hubungkan eksplisit ke Algoritma minggu ini.}
\end{frame}

\Bagian{2}{Coba dulu, baru dijelaskan}{25 menit eksperimen di GDB Online}
\note{Struggle 25 menit. Berkeliling, jangan menjelaskan materi dulu.}

\begin{frame}{Misi kalian}
\framesubtitle{Struggle, 25 menit}
\begin{columns}[T]
\begin{column}{0.31\textwidth}
\begin{Langkah}{1}{Tukar}
Tukar isi \texttt{a} dan \texttt{b} sehingga \texttt{a = 7, b = 3} menjadi \texttt{a = 3, b = 7}.
\end{Langkah}
\end{column}
\begin{column}{0.31\textwidth}
\begin{Langkah}{2}{Satu putaran}
Pada \texttt{\{5, 1, 4, 2\}}, tukar setiap pasangan bersebelahan yang terbalik dari kiri ke kanan.
\end{Langkah}
\end{column}
\begin{column}{0.31\textwidth}
\begin{Langkah}{3}{Semua putaran}
Ulangi putaran itu sampai data terurut. Berapa putaran yang dibutuhkan?
\end{Langkah}
\end{column}
\end{columns}
\vspace{8pt}
\begin{columns}[T]
\begin{column}{0.47\textwidth}
\begin{Blok}{Boleh}
Bertanya ke teman, mengubah apa saja, sengaja membuat error.
\end{Blok}
\end{column}
\begin{column}{0.47\textwidth}
\begin{BlokAlert}{Belum dulu}
Mencari jawaban jadi di internet atau AI.
\end{BlokAlert}
\end{column}
\end{columns}
\note{Catat kesalahan yang sering muncul untuk dibahas di debrief.}
\end{frame}

\begin{frame}{Apa yang kalian temukan?}
\framesubtitle{Debrief struggle}
\begin{columns}[T]
\begin{column}{0.31\textwidth}
\begin{BlokGelap}{Pertanyaan 1}
Apa yang terjadi bila menukar dengan \texttt{a = b; b = a;}?
\end{BlokGelap}
\end{column}
\begin{column}{0.31\textwidth}
\begin{BlokGelap}{Pertanyaan 2}
Sesudah satu putaran, elemen mana yang pasti sudah di tempatnya?
\end{BlokGelap}
\end{column}
\begin{column}{0.31\textwidth}
\begin{BlokGelap}{Pertanyaan 3}
Apakah selalu butuh n − 1 putaran?
\end{BlokGelap}
\end{column}
\end{columns}
\vspace{10pt}
Jawaban kalian akan kita uji di bagian berikutnya.\note{Tulis jawaban di papan; buktikan atau patahkan di bagian materi.}
\end{frame}

\Bagian{3}{Mengurutkan langkah demi langkah}{Pertukaran, bubble sort, selection sort, dan array sejajar}
\note{Materi 40 menit. Tunjukkan setiap contoh langsung di GDB Online.}

\begin{frame}[fragile]{Menukar dua nilai}
\framesubtitle{Variabel sementara}
\begin{columns}[T]
\begin{column}{0.55\textwidth}
\begin{Kode}[]{tukar.cpp}
int a = 7, b = 3;
int t = a;
a = b;
b = t;
cout << a << " " << b << endl;
\end{Kode}
\end{column}
\begin{column}{0.41\textwidth}
\only<1>{\begin{TebakDulu}
{\ITERAKickerFont\fontsize{10pt}{12pt}\selectfont\color{ITERAAlert}TEBAK DULU}\par\vspace{4pt}
Sebelum dijalankan, tulis keluarannya di kertas.
\end{TebakDulu}
}\only<2>{\KeluaranFile[]{keluaran/P12/k001.txt}
\begin{Catatan}
Tanpa \texttt{t}, nilai lama \texttt{a} hilang begitu ditimpa.
\end{Catatan}
}
\end{column}
\end{columns}
\note<1>{Minta mahasiswa menebak keluaran sebelum membuka slide berikutnya. Tanyakan satu atau dua tebakan yang berbeda, lalu buktikan.}
\end{frame}

\begin{frame}{Bubble sort: satu putaran}
\framesubtitle{Pasangan bersebelahan}
Putaran pertama pada \{5, 1, 4, 2, 8\}: bandingkan setiap pasangan, tukar bila terbalik.
\vspace{8pt}

\small
\begin{tabular}{@{}>{\raggedright\arraybackslash}p{172pt}>{\raggedright\arraybackslash}p{326pt}>{\raggedright\arraybackslash}p{364pt}@{}}
\kepala{Pasangan} & \kepala{Isi array sesudahnya} & \kepala{Keterangan}\\
a[0], a[1] & 1 5 4 2 8 & 5 > 1, tukar\\
\barisgenap a[1], a[2] & 1 4 5 2 8 & 5 > 4, tukar\\
a[2], a[3] & 1 4 2 5 8 & 5 > 2, tukar\\
\barisgenap a[3], a[4] & 1 4 2 5 8 & 5 < 8, tetap\\
\garisdasar
\end{tabular}
\vspace{10pt}

Sesudah putaran pertama, elemen terbesar pasti sudah di ujung kanan, seperti gelembung yang naik ke permukaan.
\end{frame}

\begin{frame}[fragile]{Bubble sort di C++}
\framesubtitle{Dua perulangan bersarang}
\begin{columns}[T]
\begin{column}{0.60\textwidth}
\begin{Kode}[listing options={style=ITERACodeKecil}]{bubble.cpp}
int a[5] = {5, 1, 4, 2, 8};
int n = 5;
for (int i = 0; i < n - 1; i++) {
    for (int j = 0; j < n - 1 - i; j++) {
        if (a[j] > a[j + 1]) {
            int t = a[j]; a[j] = a[j + 1]; a[j + 1] = t;
        }
    }
}
for (int k = 0; k < n; k++) cout << a[k] << " ";
cout << endl;
\end{Kode}
\end{column}
\begin{column}{0.36\textwidth}
\only<1>{\begin{TebakDulu}
{\ITERAKickerFont\fontsize{10pt}{12pt}\selectfont\color{ITERAAlert}TEBAK DULU}\par\vspace{4pt}
Sebelum dijalankan, tulis keluarannya di kertas.
\end{TebakDulu}
}\only<2>{\KeluaranFile[listing options={style=ITERAPlainKecil}]{keluaran/P12/k002.txt}
\begin{Catatan}
Batas dalam \texttt{n - 1 - i}: setiap putaran, satu elemen di ujung sudah final. \texttt{j + 1} tidak boleh melewati n − 1.
\end{Catatan}
}
\end{column}
\end{columns}
\note<1>{Minta mahasiswa menebak keluaran sebelum membuka slide berikutnya. Tanyakan satu atau dua tebakan yang berbeda, lalu buktikan.}
\end{frame}

\begin{frame}[fragile]{Selection sort}
\framesubtitle{Pilih yang terkecil}
\begin{columns}[T]
\begin{column}{0.60\textwidth}
\begin{Kode}[listing options={style=ITERACodeKecil}]{selection.cpp}
int a[5] = {29, 10, 14, 37, 13};
for (int i = 0; i < 4; i++) {
    int iMin = i;
    for (int j = i + 1; j < 5; j++) {
        if (a[j] < a[iMin]) iMin = j;
    }
    int t = a[i]; a[i] = a[iMin]; a[iMin] = t;
}
for (int k = 0; k < 5; k++) cout << a[k] << " ";
cout << endl;
\end{Kode}
\end{column}
\begin{column}{0.36\textwidth}
\only<1>{\begin{TebakDulu}
{\ITERAKickerFont\fontsize{10pt}{12pt}\selectfont\color{ITERAAlert}TEBAK DULU}\par\vspace{4pt}
Sebelum dijalankan, tulis keluarannya di kertas.
\end{TebakDulu}
}\only<2>{\KeluaranFile[listing options={style=ITERAPlainKecil}]{keluaran/P12/k003.txt}
\begin{Catatan}
Paling banyak satu pertukaran per putaran. Untuk urutan turun, cari yang terbesar.
\end{Catatan}
}
\end{column}
\end{columns}
\note<1>{Minta mahasiswa menebak keluaran sebelum membuka slide berikutnya. Tanyakan satu atau dua tebakan yang berbeda, lalu buktikan.}
\end{frame}

\begin{frame}[fragile]{Mengurutkan dua array sejajar}
\framesubtitle{Nama ikut skornya}
\begin{columns}[T]
\begin{column}{0.60\textwidth}
\begin{Kode}[listing options={style=ITERACodeKecil}]{sejajar.cpp}
string nama[4] = {"Rani", "Budi", "Siti", "Andi"};
int skor[4] = {82, 75, 95, 88};
for (int i = 0; i < 3; i++) {
    int iMaks = i;
    for (int j = i + 1; j < 4; j++)
        if (skor[j] > skor[iMaks]) iMaks = j;
    swap(skor[i], skor[iMaks]);
    swap(nama[i], nama[iMaks]);
}
for (int i = 0; i < 4; i++)
    cout << nama[i] << " " << skor[i] << endl;
\end{Kode}
\end{column}
\begin{column}{0.36\textwidth}
\only<1>{\begin{TebakDulu}
{\ITERAKickerFont\fontsize{10pt}{12pt}\selectfont\color{ITERAAlert}TEBAK DULU}\par\vspace{4pt}
Sebelum dijalankan, tulis keluarannya di kertas.
\end{TebakDulu}
}\only<2>{\KeluaranFile[listing options={style=ITERAPlainKecil}]{keluaran/P12/k004.txt}
\begin{Catatan}
Setiap pertukaran pada array kunci harus diikuti pertukaran yang sama pada array pasangannya. \texttt{swap} dari pustaka standar menukar dua nilai.
\end{Catatan}
}
\end{column}
\end{columns}
\note<1>{Minta mahasiswa menebak keluaran sebelum membuka slide berikutnya. Tanyakan satu atau dua tebakan yang berbeda, lalu buktikan.}
\end{frame}

\begin{frame}[fragile]{Tebak keluarannya}
\framesubtitle{Latihan menjelaskan, 3 menit}
\begin{columns}[T]
\begin{column}{0.56\textwidth}
\begin{Kode}[listing options={style=ITERACodeKecil}]{tebak.cpp}
int a[4] = {4, 3, 2, 1};
int tukar = 0;
for (int i = 0; i < 1; i++) {
    for (int j = 0; j < 3 - i; j++) {
        if (a[j] > a[j + 1]) {
            swap(a[j], a[j + 1]);
            tukar++;
        }
    }
}
for (int k = 0; k < 4; k++) cout << a[k];
cout << " " << tukar << endl;
\end{Kode}
\end{column}
\begin{column}{0.40\textwidth}
\begin{Catatan}
Tulis keluarannya di kertas, karakter demi karakter. Jangan dijalankan dulu.
\end{Catatan}
\end{column}
\end{columns}
\note{Contoh soal Bagian B. Beri waktu 3 menit sebelum slide jawaban.}
\end{frame}

\begin{frame}[fragile]{Tebak keluarannya: jawaban}
\framesubtitle{Latihan menjelaskan}
\begin{columns}[T]
\begin{column}{0.56\textwidth}
\begin{Kode}[listing options={style=ITERACodeKecil}]{tebak.cpp}
int a[4] = {4, 3, 2, 1};
int tukar = 0;
for (int i = 0; i < 1; i++) {
    for (int j = 0; j < 3 - i; j++) {
        if (a[j] > a[j + 1]) {
            swap(a[j], a[j + 1]);
            tukar++;
        }
    }
}
for (int k = 0; k < 4; k++) cout << a[k];
cout << " " << tukar << endl;
\end{Kode}
\end{column}
\begin{column}{0.40\textwidth}
\begin{Keluaran}[]
3214 3
\end{Keluaran}
\begin{Catatan}
Hanya satu putaran (i = 0). Angka 4 terdorong ke ujung lewat tiga pertukaran, tetapi sisanya belum terurut: 3 2 1 4.
\end{Catatan}
\end{column}
\end{columns}
\note{Tanyakan jawaban yang paling sering salah, lalu telusuri bersama.}
\end{frame}

\begin{frame}{Error yang sering muncul minggu ini}
\framesubtitle{Pesan diambil langsung dari g++}
\footnotesize
\begin{tabular}{@{}>{\raggedright\arraybackslash}p{208pt}>{\raggedright\arraybackslash}p{256pt}>{\raggedright\arraybackslash}p{398pt}@{}}
\kepala{Kesalahan} & \kepala{Contoh} & \kepala{Pesan pertama dari g++}\\
swap tanpa nilai yang bisa diubah & \texttt{swap(a[0], 5)} & \texttt{error: no matching function for call to 'swap(int\&, int)'}\\
\barisgenap Tukar string dengan int & \texttt{swap(nama[i], skor[i])} & \texttt{error: no matching function for call to 'swap(std::string\&, int\&)'}\\
Variabel sementara tak dipakai & \texttt{int t; tidak dipakai} & \texttt{warning: unused variable 't'}\\
\barisgenap Indeks bukan bilangan bulat & \texttt{a[j + 0.5]} & \texttt{error: invalid types 'int [3][double]' for array subscript}\\
Ukuran array tidak cocok & \texttt{int a[3] = \{1, 2, 3, 4\}} & \texttt{error: too many initializers for 'int [3]'}\\
\garisdasar
\end{tabular}
\vspace{10pt}

{\small Pesan di tabel ini dihasilkan dengan g++ dan opsi -Wall, sama seperti di cek.sh.}
\note{Minta mahasiswa memotret slide ini.}
\end{frame}

\Bagian{4}{Hands-on bertingkat}{45 menit, dari dasar sampai tantangan}
\note{Mahasiswa membuka folder 02 Hands-on/P12 - Sorting - Hands-on.}

\begin{frame}{Peta hands-on}
\framesubtitle{Folder 02 Hands-on/P12 - Sorting - Hands-on}
\small
\begin{tabular}{@{}>{\raggedright\arraybackslash}p{36pt}>{\raggedright\arraybackslash}p{267pt}>{\raggedright\arraybackslash}p{110pt}>{\raggedright\arraybackslash}p{83pt}>{\raggedright\arraybackslash}p{341pt}@{}}
\kepala{No} & \kepala{File} & \kepala{Tingkat} & \kepala{Waktu} & \kepala{Yang dilatih}\\
1 & \texttt{handson1-bubble.cpp} & Dasar & 10' & bubble sort naik dengan jejak per putaran\\
\barisgenap 2 & \texttt{handson2-peringkat.cpp} & Menengah & 15' & selection sort turun, dua array sejajar\\
3 & \texttt{handson3-tukar.cpp} & Tantangan & 20' & telusuri pertukaran, perbaiki tiga kesalahan\\
\barisgenap + & \texttt{bonus-hitung.cpp} & Pengayaan & sisa & bubble sort berhenti dini, hitung kerja\\
\garisdasar
\end{tabular}
\vspace{10pt}

Kerjakan berurutan. Cocokkan keluaran dengan folder \texttt{expected/}; latihan yang membaca masukan memakai data di folder \texttt{input/}.\note{Saat berkeliling, minta mahasiswa yang mengerjakan latihan tantangan menjelaskan blok PENJELASAN mereka.}
\end{frame}

\begin{frame}{Cara mengecek dan aturannya}
\framesubtitle{Hands-on}
\begin{columns}[T]
\begin{column}{0.47\textwidth}
\begin{Blok}{Di GDB Online}
Salin isi file latihan, lengkapi TODO, klik Run. Bila program membaca masukan, ketik isi file di \texttt{input/}. Bandingkan dengan \texttt{expected/}.
\end{Blok}
\end{column}
\begin{column}{0.47\textwidth}
\begin{BlokGelap}{Di komputer sendiri}
Jalankan \texttt{bash cek.sh}. Skrip mengompilasi dengan \texttt{-Wall -Werror}, memberi masukan dari \texttt{input/}, lalu mencocokkan keluaran.
\end{BlokGelap}
\end{column}
\end{columns}
\vspace{6pt}
\begin{Catatan}
AI boleh untuk memahami pesan error, tidak untuk meminta solusi utuh. Exit Ticket dikerjakan tanpa bantuan apa pun.
\end{Catatan}
\note{Hands-on tidak dinilai langsung; yang dinilai Exit Ticket dan tugas.}
\end{frame}

\Bagian{5}{Exit Ticket dan tugas}{25 menit terakhir, lalu pekerjaan rumah}
\note{Mulai Exit Ticket tepat waktu.}

\begin{frame}{Exit Ticket 12}
\framesubtitle{Individual, 25 menit, tanpa AI dan internet}
\small
\begin{tabular}{@{}>{\raggedright\arraybackslash}p{60pt}>{\raggedright\arraybackslash}p{560pt}>{\raggedright\arraybackslash}p{80pt}>{\raggedright\arraybackslash}p{150pt}@{}}
\kepala{Soal} & \kepala{Isi} & \kepala{Poin} & \kepala{CPMK}\\
A1 & Tulis bubble sort yang mengurutkan array turun & 25 & CPMK0607\\
\barisgenap A2 & Urutkan dua array sejajar (nama dan nilai) berdasarkan nilai & 25 & CPMK0607\\
B1 & Tuliskan isi array sesudah setiap putaran selection sort & 20 & CPMK0608\\
\barisgenap B2 & Jelaskan kesalahan pada potongan kode pertukaran & 20 & CPMK0608\\
B3 & Bandingkan banyak pertukaran bubble dan selection sort untuk data yang diberikan & 10 & CPMK0608\\
\garisdasar
\end{tabular}
\vspace{10pt}

{\small Soal A dinilai dari keluaran yang tepat sama. Soal B dinilai dari ketepatan dan alasan. Pelanggaran aturan membuat nilai Exit Ticket ini 0.}
\note{Soal lengkap dibagikan terpisah dan tidak disimpan di repo publik.}
\end{frame}

\begin{frame}{Tugas 12: Sistem Leaderboard}
\framesubtitle{Dikumpulkan sebelum Pertemuan 13}
\begin{columns}[T]
\begin{column}{0.47\textwidth}
\begin{Blok}{Bagian A, program, 50 poin}
Buat program leaderboard dengan menu berikut untuk maksimal 10 data (nama satu kata dan nilai). Ketentuannya di slide berikut.
\end{Blok}
\end{column}
\begin{column}{0.47\textwidth}
\begin{BlokGelap}{Bagian B, penjelasan, 50 poin}
Komentar maksimal 150 kata dan tabel isi array sesudah setiap putaran selection sort untuk lima data contoh. Jelaskan mengapa nama harus ikut ditukar, dan mengapa binary search hanya dipakai sesudah pengurutan abjad.
\end{BlokGelap}
\end{column}
\end{columns}
\vspace{8pt}
{\small Data di tugas adalah data kalian sendiri. Rubrik lengkap ada di Modul Belajar 12.}
\note{Ingatkan bahwa penjelasan disalin dari AI mudah dikenali karena tidak menyebut pengalaman mereka sendiri.}
\end{frame}

\begin{frame}{Ketentuan Tugas 12}
\framesubtitle{Sistem Leaderboard}
\footnotesize
\begin{tabular}{@{}>{\raggedright\arraybackslash}p{87pt}>{\raggedright\arraybackslash}p{788pt}@{}}
\kepala{No} & \kepala{Ketentuan Bagian A}\\
1 & Tambah data peserta dengan validasi nilai 0 sampai 100\\
\barisgenap 2 & Urutkan berdasarkan nilai turun dengan selection sort, lalu tampilkan leaderboard\\
3 & Urutkan berdasarkan nama (abjad) dengan bubble sort\\
\barisgenap 4 & Cari nama dengan linear search; tampilkan peringkat dan nilainya\\
5 & Cari nama dengan binary search pada data yang sudah urut abjad\\
\barisgenap 6 & Tampilkan banyak perbandingan dan pertukaran setiap algoritma yang dipakai, lalu keluar\\
Bonus & Bandingkan banyak langkah linear dan binary search untuk data yang sama, dan tampilkan dalam satu tabel.\\
\garisdasar
\end{tabular}
\note{Bacakan ketentuan satu per satu dan tanyakan apakah ada yang belum jelas.}
\end{frame}

\begin{frame}{Rubrik Tugas 12}
\framesubtitle{Empat level, sama di semua instrumen}
\footnotesize
\begin{tabular}{@{}>{\raggedright\arraybackslash}p{166pt}>{\raggedright\arraybackslash}p{44pt}>{\raggedright\arraybackslash}p{166pt}>{\raggedright\arraybackslash}p{156pt}>{\raggedright\arraybackslash}p{146pt}>{\raggedright\arraybackslash}p{146pt}@{}}
\kepala{Kriteria} & \kepala{Poin} & \kepala{Sangat baik 85--100} & \kepala{Baik 70--84} & \kepala{Cukup 55--69} & \kepala{Kurang <55}\\
A1 Program berjalan & 20 & tanpa error dan peringatan & ada peringatan & perlu satu perbaikan & tidak terkompilasi\\
\barisgenap A2 Kelengkapan menu & 15 & dua sort, dua search, penghitung & satu fitur kurang & dua kurang & tiga atau lebih kurang\\
A3 Ketepatan urutan dan pencarian & 15 & semua tepat termasuk nilai sama & satu kasus salah & dua salah & sebagian besar salah\\
\barisgenap B1 Jejak putaran dan alasan & 30 & tabel tepat dan alasan jelas & satu putaran keliru & tanpa alasan & tidak ada atau disalin\\
B2 Cerita error dan pengujian & 20 & pesan, sebab, perbaikan, uji & pesan tidak dikutip & hanya perbaikan & tidak ada\\
\garisdasar
\end{tabular}
\vspace{8pt}

{\small Nilai kriteria = poin × skor level. Bagian A total 50, Bagian B total 50.}
\note{Rubrik ini sama strukturnya setiap minggu; yang berubah hanya isi kriteria A2, A3, dan B1.}
\end{frame}

\begin{frame}{Sebelum Pertemuan 13}
\framesubtitle{Persiapan}
\begin{columns}[T]
\begin{column}{0.31\textwidth}
\begin{Langkah}{1}{Selesaikan}
Hands-on yang belum lulus \texttt{cek.sh}, terutama latihan tantangan.
\end{Langkah}
\end{column}
\begin{column}{0.31\textwidth}
\begin{Langkah}{2}{Kumpulkan}
Tugas 12 sebelum Pertemuan 13 dimulai.
\end{Langkah}
\end{column}
\begin{column}{0.31\textwidth}
\begin{Langkah}{3}{Baca}
Modul Belajar 12 untuk mengulang, lalu bagian awal modul berikutnya.
\end{Langkah}
\end{column}
\end{columns}
\vspace{10pt}
Pertemuan 13 membahas fungsi: memecah program menjadi bagian bernama yang bisa dipanggil berulang, termasuk fungsi pengurutan dan pencarian.\note{Tanyakan satu hal yang paling membingungkan hari ini untuk dibuka di review minggu depan.}
\end{frame}

\SlidePenutup{Terima kasih}{Sampai jumpa di Pertemuan 13: fungsi dan modularitas}
\note{Slide penutup.}
\end{document}
